% -*-memoria.tex-*-
% Este fichero es parte de la plantilla LaTeX para
% la realización de Proyectos Final de Carrera, protejido
% bajo los términos de la licencia GFDL.
% Para más información, la licencia completa viene incluida en el
% fichero fdl-1.3.tex

% Copyright (C) 2009 Pablo Recio Quijano 

%-------------------------------------------------------
% ---- Plantilla para libros / memorias PFC -----

% Realizada por Pablo Recio Quijano y Noelia Sales Montes 
% Formato de portada y primera página tomado del PFC de
% Francisco Javier Vázquez Púa, en su proyecto 'libgann'
% -------------------------------------------------------

\documentclass[a4paper,11pt]{book}

\usepackage{./estilos/estiloBase} % Basicamente son todas las
                                  % librerias usadas. En caso de que
                                  % falten librerias se van añadiendo
                                  % al fichero.
\usepackage{./estilos/colores}  % Algunos colores ya generados, para
                                % los algunos estilos más avanzados.
\usepackage{./estilos/comandos} % Algunos comandos personalizados

\graphicspath{{./imagenes/}} % Indicamos la ruta donde se encuentran
                             % las imagenes, para ahorrarnos la ruta
                             % completa, y solo modificar aquí si en
                             % un momento dado lo movemos

\begin{document}

% Renombramos las figuras y las tablas
\renewcommand{\figurename}{Figura}
\renewcommand{\listfigurename}{Índice de figuras}
\renewcommand{\tablename}{Tabla}
\renewcommand{\listtablename}{Índice de tablas}

\pagestyle{empty}
% -*-portada.tex-*-
% Este fichero es parte de la plantilla LaTeX para
% la realización de Proyectos Final de Carrera, protejido
% bajo los términos de la licencia GFDL.
% Para más información, la licencia completa viene incluida en el
% fichero fdl-1.3.tex

% Fuente tomada del PFC 'libgann' de Javier Vázquez Púa

\begin{titlepage}

  \begin{center}

    \includegraphics[scale=0.2]{logo_uca.png} \\
    
    \vspace{2.0cm}
    
    \LARGE{\textbf{ESCUELA SUPERIOR DE INGENIERÍA}} \\
    
    \vspace{1.0cm}
    
    \Large{\textbf{INGENIERÍA TÉCNICA EN INFORMÁTICA DE GESTIÓN}} \\
    
    \vspace{3.0cm}
    
    \Large{Simulador de Fútbol Online desarrollado bajo Ruby On Rails} \\
    
    \vspace{2.0cm}
    
    \Large{Carlos Rafael Belizón Ibáñez} \\
  
    \vspace{0.5cm}

    \large{\today}
    
  \end{center}
\end{titlepage}

\cleardoublepage

% -*-primerahoja.tex-*-
% Este fichero es parte de la plantilla LaTeX para
% la realización de Proyectos Final de Carrera, protejido
% bajo los términos de la licencia GFDL.
% Para más información, la licencia completa viene incluida en el
% fichero fdl-1.3.tex

% Fuente tomada del PFC 'libgann' de Javier Vázquez Púa

\begin{center}

  \includegraphics[scale=0.2]{logo_uca.png} \\

  \vspace{2.0cm}

  \Large{ESCUELA SUPERIOR DE INGENIERÍA} \\

  \vspace{1.0cm}

  \large{INGENIERO TÉCNICO EN INFORMÁTICA DE GESTIÓN} \\

  \vspace{2.0cm}

  \large{SIMULADOR DE FÚTBOL ONLINE DESARROLLADO BAJO RUBY ON RAILS} \\

  \vspace{1.0cm}

\end{center}

\begin{itemize}
\item \large{Departamento: Departamento de Ingeniería Informática}
\item \large{Directores del proyecto: Juan Manuel Dodero Beardo y Manuel Palomo Duarte}
\item \large{Autor del proyecto: Carlos Rafael Belizón Ibáñez}
\end{itemize}

\vspace{1.0cm}

\begin{flushright}
  \large{Cádiz, \today} \\

  \vspace{2.5cm}

  \large{Fdo: Carlos Rafael Belizón Ibáñez}
\end{flushright}

\cleardoublepage
\pagestyle{plain}

\frontmatter % Introducción, índices ...

% -*-previo.tex-*-
% Este fichero es parte de la plantilla LaTeX para
% la realización de Proyectos Final de Carrera, protejido
% bajo los términos de la licencia GFDL.
% Para más información, la licencia completa viene incluida en el
% fichero fdl-1.3.tex

% Copyright (C) 2009 Pablo Recio Quijano 

\section*{Agradecimientos}

Me gustaria agradecer y dedicar este texto a mis padres y hermano por
la paciencia que han tenido.

\cleardoublepage

\section*{Licencia} % Por ejemplo GFDL, aunque puede ser cualquiera

Este documento ha sido liberado bajo Licencia GFDL 1.3 (GNU Free
Documentation License). Se incluyen los términos de la licencia en
inglés al final del mismo.\\

Copyright (c) 2009 Carlos Rafael Belizón Ibáñez.\\

Permission is granted to copy, distribute and/or modify this document under the
terms of the GNU Free Documentation License, Version 1.3 or any later version
published by the Free Software Foundation; with no Invariant Sections, no
Front-Cover Texts, and no Back-Cover Texts. A copy of the license is included in
the section entitled "GNU Free Documentation License".\\

\cleardoublepage
\cleardoublepage

\tableofcontents
\listoffigures
\listoftables

\mainmatter % Contenido en si ...

\chapter{Introducción}
\label{cap:introduccion}
% -*-cap1.tex-*- Este fichero es parte de la plantilla LaTeX para la realización
% de Proyectos Final de Carrera, protejido bajo los términos de la licencia
% GFDL.  Para más información, la licencia completa viene incluida en el fichero
% fdl-1.3.tex

% Copyright (C) 2009 Pablo Recio Quijano

\section{Introducción}

Este documento describe el proceso que hemos seguido para la realización del
proyecto.  Comenzamos con la recopilación de requisitos obteniendo al tinal un
documento de Especifacación de Requisitos del Software, realizada a través de
las sucesivas entrevistas con los tutores del proyecto de fin de carrera.

Posteriormente seleccionamos la tecnología más adecuada para llevar a cabo el
proyecto. A continuación pasamos a desarrollar los componentes del software que
añadiremos al sistema para que cumpla con los requisitos que hemos recopilado
previamente.

Para finalizar, seguiremos un desarollo software ágil, por el cual vamos
desarollando partes del sistema de forma que sean totalmente funcionales una vez
terminadas hasta conformar el proyecto según las especificaciones que se han
recogido.

\section{Objetivos}

Los objetivos de este proyecto son por una parte poner en práctica los
conocimientos adquiridos a lo largo de la formación académica recibida en la
titulación de \cursiva{Ingeniería Técnica en Informática de Gestión} y a la vez
aprender a desarrollar bajo una novedosa tecnología web basada en un
\programa{framework} basado en el patrón Modelo-Vista-Controlador (\sigla{MVC})
denominado \programa{Ruby On Rails} (\sigla{MVC}).

Este \programa{framework}, a su vez, basa su simplicidad en el uso de un
lenguaje de \programa{script} llamado \programa{Ruby} que tiene características
propias de lenguajes funcionales; además, es uno de los exponentes de una
perspectiva de desarrollo en auge que se engloba dentro de las denominadas
metodologías ágiles.

El resultado de este proyecto es el desarrollo de una aplicación informática
\programa{web} de código abierto para realizar la simulación virtual de la
\cursiva{gestión computerizada} de un club de fútbol, la simulación de los
partidos que se jueguen, además de la interacción con otros usuarios. En
definitiva: un videojuego multijugador con interfaz web.

Pasamos a enumerar los objetivos de forma esquemática:

\begin{itemize}
\item \negrita{Gestión de usuarios:} todas las herramientas necesarias para
  identificar a cada usuario de nuestra aplicación web de forma unívoca.
\item \negrita{Gestión de clubes: } todas las herramientas necesarias para
  gestionar precios de entrada, venta, compra, entrenamiento de futbolistas y
  tipos de entrenamientos.
\item \negrita{Gestión de campeonatos: } todas las herramientas necesarias para
  que el administrador pueda controlar el calendario de competición, y la
  creación es de nuevas ligas en función de los usuarios que se vayan
  registrando en la web.
\item \negrita{Simulación de partidos: } desarrollaremos un simulador que
  mediante diferentes ponderaciones sea capaz de simular la realidad de un
  partido de fútbol 11.
\end{itemize}

\section{Alcance}
Se desea informatizar toda la gestión y simulación de un entorno de competición
de fútbol 11 de forma que acabemos obteniendo un videojuego multijugador
interactivo con interfaz web.

\subsection{Identificación del producto mediante un nombre}
El producto realizado se llama ``Football Simulator''.

A lo largo de toda la memoria nos referiremos a él como
\negrita{\cursiva{Simulador de fútbol online}}.

\subsection{Funcionalidades del producto}
El producto que deseamos obtener tal y como hemos mencionado anteriormente es
capaz de simular la gestión, interacción y desarrollo de una competición de liga
de fútbol 11, proporcionando funcionalidades de interacción entre usuarios y del
administrador para controlar el comportamiento de la competición.

Por otra parte, nuestro producto también es capaz de simular mediante diversas
ponderaciones los resultados de un partido de fútbol.

Para disfrutar de las funcionalidades de nuestro producto es necesario el
registro y participación de varios usuarios en él, además de la interacción de
un administrador que llevará a cabo acciones de control, supervisión y
seguimiento del desarrollo de la competición.

\subsection{Aplicaciones del software: beneficios, objetivos y metas}

Con este producto podemos obtener diversas ventajas ya descritas anteriormente,
entre ellas cabe destacar la gestión de los equipos de fútbol y la interacción
de usuarios.

El principal objetivo de la utilización del producto es divertir y entretener al
usuario, mediante la inmersión en un entorno virtual de gestión de clubes con
diversos usuarios que compiten entre sí para alcanzar la posición más alta en
clasificación y división que puedan.

Al fin y al cabo lo que se desea es formar la comunidad más grande posible y
para ello es necesario hacer que el usuario desee acceder asiduamente a nuestra
aplicación.

La mayor ventaja es que al ser una aplicación web, cualquier usuario puede
registrarse sin ningún tipo de restricción.

\section{Definiciones, acrónimos y abreviaturas}

\begin{itemize}
\item \negrita{Framework: } Estructura conceptual y tecnológica de soporte
  definida, normalmente con artefactos o módulos de sofware concretos, con base
  en la cual otro proyecto software puede ser organizado y desarrollado.
\item \negrita{Patrón Modelo Vista Controlador (MVC): } Es un estilo de
  arquitectura software que separa los datos de una aplicación, la interfaz de
  usuario y la lógica de control en tres componentes distintos.
\item \negrita{COC (Convenction over Configuration): } En español significa
  \cursiva{``Convención sobre configuración''}, quiere decir que toda pieza de
  código tiene su lugar y todas ellas interactúan según un camino estándar.
\item \negrita{DRY (Don't Repeat Yourself): } En español significa \cursiva{``No
    te repitas''}. Se refiera a que es innecesario repetir código si éste está
  bien diseñado y refactorizado.
\item \negrita{Programación Orientada a Objetos (POO): } . Es un paradigma de la
  programación que se basa en la utilización de clases de objetos.
\item \negrita{Simulador de Fútbol Online:} es el nombre utilizado comúnmente
  durante todo el documento para hacer referencia al nombre del proyecto.
\item \negrita{Ruby: } Lenguaje de programación orientado a objetos de script
  con características propias de lenguajes funcionales.
\item \negrita{Gema: }
\item \negrita{Ruby On Rails (RoR): } Framework de desarrollo rápido de
  aplicaciones web basado en el lenguaje de programación \programa{Ruby}.
\item \negrita{GIT: } Sistema de control de la configuración distribuido.
\item \negrita{Integrated Development Environment (IDE): } En español significa
  \cursiva{``Entorno integrado de desarrollo''}. Es un programa informático
  compuesto por un conjunto de herramientas y módulos de programación que ayudan
  al desarrollo de productos software.
\item \negrita{NetBeans: } IDE con funcionalidades que ayudan al desarrollo de
  aplicaciones \programa{Ruby on Rails}
\item \negrita{Sistema de Gestión de Base de Datos (SGBD)}:
\item \negrita{Oracle: } Es un sistema de Gestión de Base de Datos privativo.
\item \negrita{MySQL: } Sistema de gestión de base de datos de software libre.
\item \negrita{SQLite2: } Sistema de gestión de base de datos de software libre.
\item \negrita{SQLite3: } Sistema de gestión de base de datos de software libre.
\item \negrita{Postgre: } Sistema de gestión de base de datos de software libre.
\item \negrita{Cascade Style Sheet (CSS):} En español significa \cursiva{``Hojas
    de estilo en cascada''}.
\item \negrita{Hypertext Markup Language (HTML))} En español significa
  \cursiva{``Lenguaje de marcado de hipertexto''}.
\item \negrita{JavaScript: } Lenguaje de script que se ejecuta en el lado
  cliente.
\item \negrita{JQuery: } Librería de desarrollo rápido JavaScript.
\item \negrita{Algoritmo criptográfico: } En computación y criptografía un
  algoritmo criptográfico es un algoritmo que modifica los datos de un documento
  con el objeto de alcanzar algunas características de seguridad como
  autenticación, integridad y confidencialidad.
\item \negrita{Algoritmo MD5: } En criptografía, MD5 (abreviatura de
  Message-Digest Algorithm 5, Algoritmo de Resumen del Mensaje 5) es un
  \cursiva{algoritmo criptográfico} de 128 bits ampliamente usado.
\item \negrita{Log:} Archivo de texto plano que nos indica las todas las
  acciones que se realizan en nuestra aplicación con la finalidad de conseguir
  auditar el uso de la misma.


\end{itemize}

\section{Visión General}
Esta memoria sigue las pautas descritas en el documento de referencia
\cursiva{``Recomendaciones para la documentación del Proyecto de fin de
  Carrera''} realizada por varios profesores del departamento de lenguajes y
sistemas informáticos.

Las partes que componen el documento son:

\begin{enumerate}
\item \negrita{Introducción: } Visión general. Objetivos, alcance y estructura.
\item \negrita{Planificación: } Planificación temporal del desarrollo del
  proyecto y porcentaje de esfuerzo dedicado a las diversas tareas y fases que
  lo forman.
\item \negrita{Descripción general del proyecto: } Ampliación de la visión
  global del proyecto. Perspectiva, funciones y restricciones.
\item \negrita{Análisis: } Identificación de las metas globales, perspectivas
  del cliente y recogida de información necesaria para llevar a cabo el
  proyecto.
\item \negrita{Diseño: } Aplicamos diferentes técnicas y procedimientos para
  obtener un producto con el suficiente detalle que permita su realización
  física.
\item \negrita{Implementación: } Aspectos más relevantes de la implementación de
  la aplicación software.
\item \negrita{Pruebas: } Todo lo concerniente a las pruebas para que la
  aplicación logre su cometido sin fallos.
\item \negrita{Resumen: } Breve resumen de lo más destacable del proyecto
  realizado.
\item \negrita{Conclusiones y trabajo futuro: } Valoración global, posibles mejoras y
  ampliaciones del proyecto.
\item \negrita{Apéndices: }
  \begin{itemize}
  \item \negrita{Manual de usuario: } Manual del funcionamiento de la aplicación
    destinado a los usuarios de la misma.
  \item \negrita{Manual de instalación y puesta en funcionamiento: }
    Instrucciones para realizar la instalación de la aplicación, así como su
    puesta en funcionamiento.
  \item \negrita{Difusión: } Lugares dónde se ha publicado el proyecto.
  \item \negrita{Licencia: } Texto de la licencia adoptada para la divulgación
    del proyecto.
  \end{itemize}
\item \negrita{Bibliografía: } Libros y referencias físicas y electrónicas
  consultadas para la elaboración del proyecto.
\end{enumerate}



\chapter{Planificación}
\label{cap:planificacion}
% -*-cap2.tex-*- Este fichero es parte de la plantilla LaTeX para la realización
% de Proyectos Final de Carrera, protejido bajo los términos de la licencia
% GFDL.  Para más información, la licencia completa viene incluida en el fichero
% fdl-1.3.tex

% Copyright (C) 2009 Pablo Recio Quijano

\section{Metodología de desarrollo}
La distribución temporal seguida en líneas generales es la descrita a
continuación, de la cual se harán revisiones constantes al tratarse de una
aplicación en la que se desea aprender a desarrollar bajo una metodología ágil.
\begin{itemize}
\item Análisis: En esta fase llevaremos a cabo la recogida de todos los
  requisitos funcionales de nuestra aplicación. Al utilizar una metodología ágil
  la recogida de requisitos podrá ser más laxa que si hubiéramos seguido un
  modelo lineal.
\item Diseño: En esta fase recrearemos diseños conceptuales de cómo serán
  nuestras vistas, y a partir de ellas podremos obtener nuestros modelos de
  datos.
\item Programación: En esta fase nos dedicaremos a codificar de forma sinérgica
  las diferentes secciones en las que se compone un desarrollo basado en el
  patrón Modelo Vista Controlador
\item Pruebas: Las pruebas que realizaremos serán:
  \begin{enumerate}
  \item Test de usabilidad
  \item Pruebas de caja negra.
  \item Estudio y resolución de los resultados de las pruebas.
  \end{enumerate}
\item Documentación: Para la documentación del código usaremos la herramienta
  propia del lenguaje Ruby, Ruby-doc, gracias a esta aplicación podremos generar
  un documento web fácilmente extensible y manipulable. En esta fase también
  desarrollaremos la memoria de la aplicación, el manual de usuario y la
  presentación multimedia para difundir nuestra aplicación.
\end{itemize}

\section{Calendario}

En la figura \ref{fig:gantt} se puede ver un diagrama de Gantt con la división y el tiempo
tomado para realizar cada una de las tareas en las que hemos dividido el
proyecto. La unidad de tiempo básica está expresada en días laborables de 8
horas.

Aunque la etapa de Memoria se ha realizado de forma solidaria con las demás,
para simplificar se ha extraído en un sólo bloque.

\begin{itemize}
\item Entrevistas: Nos reunimos con nuestros tutores de proyecto para recopilar
  los requisitos que deberá tener nuestro proyecto de fin de carrera para que
  sea lo suficientemente completo para pasar la aprobación del tribunal y
  aprender de forma exhausta el framework \cursiva{Ruby On Rails} (5 días).
\item Análisis: Una vez recopilada toda la información en las entrevistas con
  nuestros tutores, extraemos la información importante y la esquematizamos (7
  días).
\item Especificación de Requisitos: Definimos, organizamos, cada uno de los
  requisitos que hayamos extraído y los desarrollamos de forma exhausta. En esta
  fase también describiremos los casos de uso, evitando en la medida de lo
  posible la redundancia de los mismos y maximizando la posibilidad de
  reutilizarlos (10 días).
\item Diseño: En esta fase estudiaremos la mejor forma para incorporar la
  funcionalidad del sistema teniendo en cuenta que usaremos el framework de
  desarrollo rápido de aplicaciones \cursiva{Ruby On Rails}, con especial
  atención a que haremos uso del patrón Modelo-Vista-Controlador (25 días).
\item Implementación: En esta fase codificaremos la aplicación para que sea
  funcional según hemos especificado en la fase de diseño; además, será
  necesario aprender la forma de utilizar de forma correcta el framework
  \cursiva{Ruby On Rails)}. En esta fase va incluida la creación de la
  estructuras necesarias para la interacción con el Sistema de Gestión de Bases
  de Datos (\cursiva{SGBD}) a través de un sistema \cursiva{ORM} (90 días).
\item Pruebas: En este paso comprobaremos que la corrección del sistema que
  hemos desarrollado. En caso de encontrar algún fallo es el momento de
  solucionarlo y repetir las pruebas (5 días).
\item Memoria: Toda la documentación que conforma esta memoria es el resultado
  de detallar cada uno de los documentos que hemos ido generando a través de la
  realización de la misma (40 días).
\item Presentación: Creación y preparación de la presentación que mostraremos en
  la etapa de defensa ante el tribunal (10 días).
\end{itemize}

\section{Presupuesto}

Para realizar una estimación del montante económico del proyecto realizado será
necesario calcular y dirimir todos los recursos que han sido necesarios para
llevarlo a cabo. En dicho presupuesto no se incluirán los materiales de oficina
como serían escritorio, folios, ordenador, impresora, y demás recursos
ofimáticos. Sabiendo que excluiremos todos los gastos que podríamos considerar
como indirectos, sólo tendremos que fijarnos en el gasto que supondría el coste
de mantener a un \cursiva{Analista-programador y Diseñador de página Web} en la
empresa. Para ello consultaremos las tablas salariales de dicha categoría
profesional en el \cursiva{XVI CONVENIO COLECTIVO ESTATAL DE EMPRESAS
  CONSULTORAS, DE PLANIFICACIÓN, ORGANIZACIÓN DE EMPRESA Y CONTABLE, EMPRESAS DE
  SERVICIOS DE INFORMÁTICA Y DE ESTUDIOS DE MERCADO Y DE LA OPINIÓN PÚBLICA} tal
y como vemos en la tabla \ref{tab:salario}. Para el cálculo del pago de la
Seguridad Social aplicaremos un 33\% del coste del salario bruto del trabajador.

\begin{table}[H]
  \begin{center}
  \begin{tabular}{| c | c | c | c |}
    \hline
    Nombre & Coste mensual & Coste anual & Coste jornada (25 días por mes)\\ \hline
    Analista-programador y Diseñador Web & 1.468,54 \euro & 20.559,56
    & 58,75 \euro\\ \hline
    Seguridad Social & 484,82 \euro & 6.784,65 \euro & No aplica \\
    \hline
  \end{tabular}
\end{center}
\caption{Salario según último convenio de las TIC}
\label{tab:salario}
\end{table}


Teniendo en cuenta que el proyecto lo hemos realizado en solitario y nos ha
llevado un total de 192 jornadas laborales, lo que equivaldría a 7 meses y 3
semanas naturales, tendremos que el coste total del proyecto es el referido en
la tabla \ref{tab:coste}

\begin{table}[H]
  \begin{center}
  \begin{tabular}{| c | c | c | c |}
    \hline
    Cantidad & Descripción & Coste unitario & Coste Total\\ \hline
    192 & Jornada Analista-Programador y Diseñador Web & 58,75 \euro &
    11.280 \euro\\ \hline
    8 & Seguridad Social Analista-Programador y Diseñador Web & 484,82
    \euro & 3.878,56 \euro \\ \hline
    \multicolumn{3}{|r|}{Total} &  15.158,56 \euro\\
    \hline
  \end{tabular}
\end{center}
\caption{Coste total del proyecto realizado}
\label{tab:coste}
\end{table}

\chapter{Descripción General del proyecto}
\label{cap:descripcion}
\section{Perspectiva del Producto}
En este capítulo explicaremos qué pretendemos conseguir con la realización de
este proyecto y qué características generales ha de cumplir.

\subsection{Dependencias del producto}
El proyecto es independiente, no es un subsistema de un proyecto de mayor
envergadura, tampoco es continuación de ningún otro proyecto, pero al estar
liberado con una licencia de \cursiva{software libre} podría ser continuado,
mejorado y actualizado por su propio creador o terceras personas ajenas a éste.

\subsection{Interfaces de usuario}
La interfaz de usuario se basa en una página web dinámica que lleva asociada
varias opciones para poder usar las herramientas
proporcionadas.\\

Para el uso de la aplicación, la página principal dispone de varios menús, entre
ellos la de selección de idiomas y un \cursiva{hipervínculo} al
subsistema de registro e inicio de sesión.\\

Todo el diseño de la aplicación respeta los estándares \cursiva{HTML 4.0} y
\cursiva{CSS 2.0 } así como el uso de la librería de \cursiva{JavaScript JQuery}
lo que nos proporciona una gran
compatibilidad con cualquier navegador web del mercado.\\

Al usar \cursiva{hojas de estilo CSS} la apariencia de nuestra aplicación puede
ser fácilmente modificada para adaptarla a nuevas
tendencias de diseño de forma sencilla y eficiente.\\

Siempre que se termine cualquier operación, ya sea exitosa o no, el sistema nos
informará a través de un sistema de notificación embebido dentro del navegador
muy parecido al sistema usado en \cursiva{Mac OS
  X } y últimas versiones de la \cursiva{distribución Linux Ubuntu}.\\

Como mejora para un uso mucho más cómodo y rápido por parte de cualquier
usuario, la aplicación dispone de un menú estático para acceder a cualquier tipo
de herramienta o funcionalidad según el usuario sea un jugador de la aplicación
o sea el administrador.

\subsection{Interfaces software}
El producto interactúa con el sistema operativo en el lado del servidor en dónde
está instalado y con el navegador usado en el lado
del cliente.\\

De hecho, el producto es totalmente multiplataforma tanto en el lado servidor
como en el lado cliente, gracias a que el \cursiva{framework RoR} está diseñado
para ser instalado en entornos \cursiva{Windows, Linux y Mac OS X}; mientras que
en el lado cliente disponemos de varios navegadores web cualquiera que sea el
sistema operativo en el que se
ejecute.\\

Debemos destacar que el mencionado \cursiva{framework} tiene una gran comunidad
que lo sustenta, que da como resultado la disponibilidad de una gran cantidad de
documentación y un buen soporte de la plataforma,
repercutiendo directamente en la calidad del software que desarrollemos.\\

Hemos de indicar que \cursiva{RoR} es agnóstico en cuanto al \cursiva{SGBD} que
utilicemos, con lo que el producto podrá ser desplegado con cualquiera de los
\cursiva{SGBD} más comunes como son \cursiva{Oracle, MySQL, SQLite2, SQLite3 y
  Postgre}, a través de una \cursiva{gema} específica y una simple modificación
en uno de los archivos de configuración.

\subsection{Operaciones}
Existen dos tipos de usuarios dentro de la aplicación:

\begin{itemize}
\item \negrita{Superusuario: } Es el responsable de instalar, mantener,
  interconectar con la \cursiva{base de datos} y optimizar el servidor dónde se
  encuentra la aplicación.
\item \negrita{Administrador: } Es el responsable de toda la gestión de
  campeonatos, así como de dar de alta a los nuevos usuarios para que puedan
  empezar a interactuar con las funcionalidades de la aplicación.
\item \negrita{Usuario: } Conforma el resto de usuarios del sistema que
  participan en la simulación del entorno de gestión de equipos de fútbol, a
  excepción de las funcionalidades exclusivas del administrador.
\end{itemize}

\subsection{Requisitos de adaptación a la ubicación}
Para poder adaptar la aplicación al entorno dónde se instalará será necesario
disponer de un \cursiva{SGBD} compatible con \cursiva{RoR} además de modificar
el fichero de configuración para conectarlo con nuestro sistema.

\subsection{Funciones del producto}
La función principal del producto que estamos presentando es la de desarrollar
un videojuego multijugador dónde diversos usuarios puedan interactuar en un
entorno que simule ser real a través de \cursiva{Internet}.  Para ello ha sido
necesario crear una aplicación informática encargada de poder registrar a
cualquier usuario que acceda a nuestra web y proporcionarles las herramientas
necesarias para gestionar su propio club; además de simular los partidos que en
una competición de fútbol
se desarrollan.\\

Las características más importantes de esta aplicación se describen a
continuación:

\begin{itemize}
\item Posibilidad de registro y control de acceso a una cuenta de usuario
  personal e intransferible que distinga al usuario unívocamente de los demás a
  través de una simple interfaz web.
\item Capacidad para la gestión de todos los aspectos económicos de los clubes
  de los usuarios registrados.
\item Capacidad para la visualización del calendario y clasificación de todos
  los clubes y usuarios registrados en el sistema.
\item Capacidad para modificar la estructura deportiva y técnica de los clubes
  de los usuarios registrados.
\item Capacidad para simular un partido de forma realista y mostrarlo a los
  usuarios como si se estuviera desarrollando en directo.
\item Capacidad para gestionar el estado de la competición.
\item Capacidad para escoger el idioma en el que se mostrará la información de
  nuestra aplicación.
\end{itemize}

Aquí sólo se han descrito las características generales que realiza la
aplicación, se podrá comprobar toda la funcionalidad en los puntos siguientes de
esta memoria, en especial en el \cursiva{Manual de usuario}.

\section{Características de los usuarios}

Para el correcto manejo de la aplicación \cursiva{Simulador de fútbol
  online} existen dos perfiles distintos de manejo.\\

Podríamos dividir los perfiles en:

\begin{itemize}
\item \negrita{Usuario (Jugador):} Necesitan unos conocimientos básicos de
  informática para su manejo, sobre todo orientado al manejo de aplicaciones
  web. También será necesario un conocimiento meramente básico de cuáles son las
  reglas de negocio y deportivas de un entorno de gestión de clubes deportivos
  de fútbol 11.
\item \negrita{Usuario (Administrador): } Necesitará unos conocimientos básicos
  en informático para su manejo, sobre todo orientado al manejo de aplicaciones
  web. También será necesario conocer cómo es el discurrir de las competiciones
  en el diseño de nuestra aplicación.
\item \negrita{Superusuario (Administrador): } Serán necesarios unos
  conocimientos avanzados de informática, incidiendo en conocimientos de
  despliegue de aplicaciones web.
\end{itemize}

\section{Restricciones generales}
El proyecto se ha desarrollado bajo la metodología \cursiva{Rational Unified
  Process (RUP)}, proceso de desarrollo orientado a objetos, al estar basado
todo el proyecto en un entorno de desarrollo bajo el paradigma
de programación orientada a objetos.\\

El lenguaje de programación que usaremos será \cursiva{Ruby}, que es libre así
como el \cursiva{framework} sobre el que apoyaremos nuestra
aplicación web \cursiva{RoR}.\\

No se deben tener en cuenta restricciones en el uso de memoria principal por
parte de la aplicación tanto en el lado servidor como en
el cliente.\\

Se deben tener en cuenta restricciones a la hora de crear las vistas de nuestra
aplicación web para que éstas sean lo mas livianas posibles
para ahorrar el mayor ancho de banda posible.\\

Para el desarrollo de todo el proyecto deberemos usar un \cursiva{IDE} que se
base en la filosofía de \cursiva{software libre}, en este caso
\cursiva{NetBeans}.\\

El producto habrá de funcionar en cualquier sistema operativo que tenga
disponibilidad para tener instalado un intérprete de \cursiva{Ruby} en el lado
servidor y en cualquier navegador del mercado compatible con los estándares web
para \cursiva{HTML 4.0}, \cursiva{CSS 2.0}, y la librería \cursiva{JQuery}. Esto
hará que
nuestro software tenga un mayor público objetivo de desarrolladores.\\

Como sistema gestor de bases de datos se podrá utilizar cualquiera que tenga una
\cursiva{gema} que sea capaz de conectar \cursiva{RoR} con
el \cursiva{SGBD} escogido.\\

Es necesario tener en cuenta que nuestra aplicación puede correr en cualquier
\cursiva{pc} pero es conveniente que éste tenga la capacidad suficiente para
atender varias peticiones a la vez, además de tener la configuración necesaria
para que peticiones que lleguen desde
\cursiva{internet} puedan ser atendidas.\\

Para toda la creación y generación de la documentación presente hemos
utilizado LaTeX\\

Como controlador de la configuración hemos optado por usar \cursiva{GIT}, un
sistema de control de versiones distribuido que es el usado oficialmente por el
proyecto \cursiva{RoR} y que se integra de forma plena con \cursiva{NetBeans}.

\section{Requisitos no funcionales del sistema}
\subsection{Requisitos Generales}
La solución debe cumplir como mínimo con las siguientes características basadas
en las especificaciones funcionales y los requisitos no funcionales:

\begin{itemize}
\item Basada en la web.
\item La aplicación debe esta diseñada y desarrollada sobre la plataforma
  \cursiva{Ruby On Rails}.
\item La aplicación debe ser escalable bajo estrategia vertical (Añadir más
  recursos al servidor) y horizontal (Añadir más servidores), según las
  necesidades de procesamiento.
\item La aplicación debe tener bajo nivel de acoplamiento y la posibilidad de
  editar fácilmente los parámetros que se consideren dinámicos y requieran
  cambios frecuentes.
\item Orientada a objetos.
\item De fácil mantenimiento en cuanto a cumplimiento de estándares, uso de
  guías y patrones con especial énfasis en el patrón \cursiva{MVC},
  documentación y de fácil ubicación de componentes.
\item Que permita y utilice la reutilización de código.
\item Basada en la arquitectura cliente-servidor.
\item La aplicación debe permitir generar avisos a través de correo electrónico.
\end{itemize}

\subsection{Requisitos específicos}
Los requisitos no funcionales generales estarán enmarcados en los siguientes
aspectos:
\begin{itemize}
\item Escalabilidad:
  \begin{itemize}
  \item El diseño debe contemplar el uso óptimo de recursos como las conexiones
    con la base de datos.
  \item El diseño debe contemplar la clara división entre datos, recursos, y
    aplicaciones para optimizar la escalabilidad del sistema.
  \item Debe contemplar requisitos de crecimiento para usuarios internos al
    sistema.
  \end{itemize}
\item Disponibilidad:
  \begin{itemize}
  \item La disponibilidad del sistema ha de ser continua con un nivel de
    servicio para los usuarios de siete días y 24 horas.
  \item En caso de fallos de algún subsistema no debe haber pérdida de
    información
  \item Debe contemplar requisitos de consistencia transaccional.
  \end{itemize}
\item Seguridad:
  \begin{itemize}
  \item La aplicación debe reflejar patrones de seguridad teniendo para cumplir
    con la ley orgánica de protección de datos.
  \end{itemize}
\item Mantenibilidad:
  \begin{itemize}
  \item Se debe estructura el código de manera consistente y predecible.
  \item Aprovechar al máximo las facilidades que nos aporta el
    \cursiva{framework Ruby On Rails}, con especial énfasis en el patrón
    \cursiva{MVC}.
  \end{itemize}
\item Rendimiento: La aplicación debe ofrecer una buena respuesta ante alta
  demanda de usuarios.
\end{itemize}
\section{Requisitos para futuras versiones}
Para futuras versiones del producto, podrían introducirse funciones avanzadas
para los usuarios como selección de escudo del equipo, personalización de
dorsales, avisos por sms, mercados de agentes
libres, cantera y ojeadores.\\

Podría introducirse nuevas plantillas de diseño para personalizar la vista de
diversas formas para que el usuario escoja la que más le
guste.\\

También sería interesante establecer un apartado de estadísticas para poder
llevar un seguimiento del historial en tiempo de la evolución de los equipos
tanto en el plano técnico como en el económico, este apartado también sería
aplicable a los futbolistas que integran los
clubes.\\

Podría añadirse un sistema de foros y chat para una interacción entre usuarios
que no se limite en exclusiva a las funcionalidades
específicas de nuestro sistema.\\

Como restricción a futuras versiones o proyectos que estén basados en nuestra
aplicación se aplicarán las referentes a las propias especificadas en el
apartado de licencia de este documento.


\chapter{Análisis}
\label{cap:analisis}
% -*-cap4.tex-*- Este fichero es parte de la plantilla LaTeX para la realización
% de Proyectos Final de Carrera, protejido bajo los términos de la licencia
% GFDL.  Para más información, la licencia completa viene incluida en el fichero
% fdl-1.3.tex Copyright (C) 2009 Pablo Recio Quijano

En este capítulo vamos tratar de recopilar los requisitos funcionales de nuestro
proyecto.

\section{Metodología de desarrollo}
Hemos decidido seguir una metodología de desarrollo que se adapte lo
suficientemente bien al framework en torno al cual desarrollaremos
nuestra aplicación.\\

Hemos decidido usar la metodología de desarrollo \cursiva{Xtreme
  Programming (XP)}\\

La metodología será seguida dentro de unos límites y adaptada a nuestra
necesidades, debido entre otras cosas a que somos nosotros mismos los que hemos
decidido qué funcionalidades y qué problemas queremos resolver con nuestra
aplicación; además, no hay que olvidar que no tenemos ningún equipo que
coordinar ya que sólo será una persona quién se encargue de realizar todas y
cada una de las funcionalidades.\\

La programación extrema, o \cursiva{Xtreme Programming} es una de las llamadas
\cursiva{Metodologías Ágiles} de desarrollo de software más populares hoy en
día. Inicialmente fue descrita por \cursiva{Kent Beck} cuando trabajaba en la
\cursiva{Chrysler Corporation}. En la \cursiva{XP} se da por supuesto que no se
debe prever todo antes de empezar a codificar, que es imposible capturar todos
los requisitos del sistema, ni saber qué es todo lo que debe hacer y por tanto
es imposible hacer un diseño correcto desde el principio.\\

La idea principal de esta metodología consiste en trabajar estrechamente con el
cliente, diseñando pequeñas versiones frecuentemente. El objetivo de estas
versiones no es otro que conseguir que la aplicación funcione de la forma más
simple y eficiente posible con el mínimo código. Cuando el cliente le transmite
al programador lo que necesita, éste puede hacer una estimación aproximada del
tiempo que le llevará codificarlo, aunque lógicamente dicha planificación deberá
revisarse y modificarse continuamente a lo largo del desarrollo del proyecto, es
decir, se va iterando. Cada vez que el desarrollador consigue una versión
funcional de lo que el cliente ha solicitado, ésta se le muestra al cliente para
que la testee y haga peticiones de las modificaciones que crea convenientes. De
esta manera se evita perder tiempo desarrollando una aplicación que no sea la
que el cliente esperaba. Este ciclo se repetirá tantas veces como el cliente
necesite para sentirse satisfecho con la aplicación.\\

La \cursiva{Xtreme Programming} agrupa trece prácticas básicas que se deben
cumplir para asegurar el éxito del proyecto, ellas son:

\begin{itemize}
\item Equipo completo.
\item Planificación continua.
\item Test del cliente.
\item Versiones pequeñas.
\item Diseño simple.
\item Programación por parejas.
\item Desarrollo guiado por pruebas automáticas.
\item Mejora continua del diseño.
\item Integración continua.
\item Código en multipropiedad.
\item Normas de codificación comunes.
\item Utilización de metáforas.
\item Mantener un ritmo sostenible.
\end{itemize}

La \cursiva{Xtreme Programming} consta de unos valores, unos principios
fundamentales y unas prácticas. Los principios fundamentales en los que se basa
son:

\begin{itemize}
\item Feedback: Retroalimentación veloz.
\item Modificaciones incrementales.
\item Refactorización
\item Asunción de simplicidad.
\item Respeto: Comunicación entre desarrolladores.
\end{itemize}



\section{Descripción General}
A continuación vamos a ver los factores que afectan al producto y sus
requerimentos

\section{Propósito}
El propósito del siguiente apartado es definir cuáles son los requisitos que
debe tener nuestro portal web que simula la gestión de
un club de fútbol en un entorno multijugador y simular partidos.\\

Esta especificación de requisitos está destinada a ser leída por los usuarios o
cualquier sujeto que tenga interés en saber cómo funciona el producto.

\section{Ámbito}
El producto que vamos a describir puede clasificarse como un videojuego
multijugador online integrado en un servidor web en la cual los usuarios del
portal se dedicarán a gestionar un club de fútbol en la parcela económica y
deportiva, y dispondrá de un simulador de los
partidos entre clubes.\\

El producto ha de ofrecer una interfaz lo bastante sencilla para que cualquier
usuario pueda gestionar las distintas facetas en escasos minutos. Dado que la
aplicación es de tipo web, deberemos de tener especial cuidado en respetar los
estándares para que la aplicación sea visible en cualquier navegador del mercado
obteniendo así un videojuego que sea totalmente multiplataforma. La tecnología
que usaremos en el lado del servidor web será el framework Ruby On Rails.

\section{Perspectiva del Producto}
El presente producto ha de correr en un servidor que sea compatible con
\programa{Ruby on Rails}. Antes que nada \programa{Ruby On Rails} es un
framework que hace mas fácil desarrollar, implementar y mantener aplicaciones
web. Esto es debido entre otras razones a que:

\begin{itemize}
\item Toda aplicación web implementada usa el esquema \sigla{MVC} (Modelo Vista
  Controlador).
  
\item Toda pieza de código tiene su lugar y todas ellas interactúan según un
  camino estándar, es lo que ha venido a llamarse \sigla{COC} (Convención sobre
  Configuración).
  
\item Las aplicaciones en Rails están escritas en Ruby, un lenguaje de script
  moderno y orientado a objetos que ayuda a que se potencie otro de los pilares
  de la filosofía de desarrollo que viene abreviado en la siguiente
  nomenclatura: \sigla{DRY} (No te repitas).
\end{itemize}

En cuanto a la máquina necesaria para ejecutar el servidor, valdrá cualquier
sistema que sea compatible con \programa{Windows}, \programa{Linux} o
\programa{MacOS X} recomendándose que se use un hardware acorde con las
peticiones que el portal reciba. Los usuarios deberán de acceder al portal
mediante cualquier navegador web disponible en el mercado que sea compatible con
los estándares web \sigla{XHTML 1.1 Strict} y \sigla{CSS 2.0.}

\section{Características del usuario}
A continuación vamos a ver a qué tipo de usuarios está dirigido el producto y
cómo afectan éstos a las funciones que debe realizar
nuestro producto.\\

El producto está dirigido a todo usuario que sepa manejar un navegador web y que
esté interesado en la simulación de la gestión de un club de fútbol, así como
conocer las reglas y premisas básicas de negocio del deporte.
\section{Obligaciones generales}
Uno de los aspectos que influirá determinantemente en el éxito del producto será
la eficiencia en las peticiones y el diseño de una intertaz amigable y sencilla
para poder gestionar los diferentes
parámetros del club de una forma rápida y precisa.\\

De cara al diseño de la interfaz web, uno de los mayores desafíos con los que
nos encontraremos será el conseguir que el diseño de la interfaz sea operativa
en cualquiera de los navegadores del mercado, ya que los dos navegadores más
populares no soportan actualmente con suficiente solvencia los estándares
\sigla{CSS 2.0} (\programa{Internet Explorer 8} y \programa{Firefox 3.5}).
\section{Asunciones y Dependencias}
El presente producto depende enteramente de que el usuario disponga de una
conexión red al servidor a través de un navegador web (probablemente
\sigla{internet}), por tanto es importante respetar los
estándares pertinente que estipula la \sigla{w3c}.\\

\section{Especificación de los requisitos del sistema}
\subsection{Requisitos de interfaces externas}
\subsubsection{Requisitos de interfaces de usuario}
Podemos dividir los requisitos de interfaz de usuario según tres roles:

\begin{itemize}
\item \negrita{Usuario Jugador y Usuario Administrador: }
  \begin{itemize}
  \item Interfaz atractiva para los usuarios: Es fundamental para la buena
    aceptación de nuestra aplicación el crear una interfaz lo suficientemente
    clara, dinámica y estética para que el usuario sienta una inmersión en la
    experiencia de juego. Una interfaz no atractiva sólo traerá problemas de
    rechazo y disconformidad por parte del usuario.\\\\
    Para que la interfaz sea sencilla de modificar y tratándose de una
    aplicación web usaremos \cursiva{hojas de estilo CSS}, así todo el apartado
    de diseño gráfico podrá ser modificado posteriormente sin que por ello
    afecte al contenido ni la estructura del documento \cursiva{HTML}.
  \item Una característica de nuestra aplicación es que el sistema de avisos
    está embebido dentro de la página web y es muy parecido al sistema de
    notificación de los sistemas operativos \cursiva{Mac OS X} y de la \cursiva
    {distribución Linux Ubuntu}.
  \end{itemize}
\item \negrita{Superusuario: } En este caso la interfaz corresponderá a la que
  tenga el propio sistema operativo dónde esté instalado la aplicación.
\end{itemize}

\subsection{Requisitos de rendimiento}
Este aspecto es fundamental en el análisis del sistema, se ha de tener en cuenta
que la aplicación es de tipo servidor y que a ella estarán
conectados concurrentemente un número de usuarios elevado.\\

Se ha considerado también el rendimiento que nos ofrece el diseño web
realizado. Para que nuestra página web sea compatible con la mayoría de
dispositivos y se ejecute lo suficientemente fluida hemos decidido no usar la
tecnología propietaria \cursiva{Flash}, usando para la renderización de efectos
la librería \cursiva{JQuery} en combinación
con hojas de estilo \cursiva{CSS}.\\

El rendimiento de la aplicación podrá ser menor en la primera instanciación de
la aplicación debido a que no se habrá cacheado ningún contenido.

\subsection{Requisitos de Información}
La \cursiva{base de datos} es uno de los componentes fundamentales de la
aplicación. En ella se van a almacenar todos los datos registrados
por la aplicación.\\

La aplicación sólo hará uso de una conexión al \cursiva{SGBD} escogido por el
\cursiva{superusuario} en la configuración de instalación de la aplicación, cuyo
nombre también podrá ser escogido por el citado
\cursiva{superusuario}.\\

Dicha \cursiva{BD} estará formada por las siguientes tablas:

\begin{itemize}
\item User Tabla donde se almacenan los datos de cada uno de los usuarios que
  interactúan con la aplicación.
\item Club Tabla donde se almacenan los datos de los clubes que participan en
  las ligas.
\item League Tabla donde se almacenan los datos de las ligas que existen en el
  sistema.
\item Season Tabla donde se almacenan los datos de las temporadas
  correspondientes a cada liga del sistema.
\item Round Tabla donde se almacenan los datos de las jornadas de cada
  temporada.
\item MatchGeneral Tabla dónde se almacenan los partidos que se han de jugar
  cada jornada.
\item MatchDetail Tabla dónde se almacenan las acciones que tienen lugar en cada
  partido.
\item LineUp Tabla dónde se almacenan las alineaciones de cada equipo y partido.
\item ClubFinancesRound Tabla donde se almacenan los datos de finanzas
  correspondientes a cada jornada y club.
\item Player Tabla donde se almacenan los datos de los jugadores de fútbol de
  los clubes.
\item Offer Tabla donde se almacenan los datos de las ofertas que se realizan
  los usuarios de la aplicación sobre los jugadores de fútbol.
\item Training Tabla donde se almacenan los datos de los entrenamientos de los
  jugadores de fútbol.
\item AdminMessages Tabla donde se almacenan los mensajes de administración
  hacia los usuarios.
\end{itemize}

Toda la información sobre la \cursiva{base de datos} será ampliada en
su sección correspondiente.\\

La tabla \cursiva{User} contiene los datos de los usuarios que interactúan con
la aplicación, en ella se almacenará la contraseña de cada uno de ellos; es por
ello que será necesario un \cursiva{algoritmo criptográfico} denominado
\cursiva{MD5}. De esta forma conseguimos proteger la contraseña de los usuarios
frente a un
posible ataque al servidor dónde esté la \cursiva{base de datos}.\\

Es necesario que el \cursiva{superusuario} conozca el \cursiva{SGBD} que será
utilizado con nuestra aplicación, además de disponer de un usuario y contraseña
para la correcta creación de las tablas.

\subsection{Restricciones de diseño}

No existe ninguna restricción considerable en el diseño del sistema.\\

Todo el diseño del sistema se debe de regir con la especificación
realizada.\\

Es importante estimar cuál va ser el tráfico que nuestra aplicación va sufrir
para disponer de un \cursiva{servidor} y una conexión lo suficientemente
potentes para dar servicio a todos los usuarios que se conecten a la aplicación.

\subsection{Atributos del sistema software}
El sistema ha de cumplir las siguientes propiedades:

\begin{itemize}
\item \negrita{Fiabilidad}: Este sistema, como todo software de calidad, debe
  ser fiable. Al estar tratando de una aplicación web que ha de tener un elevado
  acceso concurrente a los datos, hemos de cerciorarnos que las operaciones no
  dejen el sistema en un
  estado inconsistente por una posible caída del servidor.\\

  Por ello el sistema se encarga de generar un \cursiva{log} con todas las
  operaciones que ocurren en nuestra aplicación.

  Hemos de tener en cuenta que el usuario de nuestra aplicación no ha de conocer
  nada sobre la estructura interna de nuestra aplicación, por ello los mensajes
  de error serán lo más simples posibles para evitar posibles ataques al
  servidor dónde es alojada.

\item \negrita{Disponibilidad}: La disponibilidad de este producto debe ser
  plena. La aplicación ha de estar disponible en cualquier momento y ser
  accesible desde cualquier lugar. Para llevar a cabo este precepto será
  necesario alojar la aplicación en un lugar web que tenga la suficiente
  capacidad para dar servicio a un elevado número de usuarios.

\item \negrita{Seguridad}: La seguridad de este tipo de aplicaciones ha de ser
  bastante elevada para conseguir cumplir leyes de protección de datos y
  asegurar un correcto funcionamiento.

  Para la seguridad y confidencialidad de los datos se debe seguir un protocolo
  dirigido a asegurar la confidencialidad de los datos almacenados en la
  \cursiva{base de datos}.

  Será necesario que el superusuario de la aplicación disponga de un nombre de
  usuario y contraseña lo suficientemente seguros para trabajar con el
  \cursiva{SGBD} escogido.

  Los usuarios de la aplicación tendrán un nombre usuario y una contraseña
  asociada que encriptada con el algoritmo \cursiva{MD5}.

  Es posible que cualquier persona acceda a la página web para estar informado
  de los que ocurre en las diferentes ligas que componen el juego. Este tipo de
  usuarios no necesitan disponer de nombre de usuario y contraseña asociada.

\item \negrita{Mantenibilidad}: La mantenibilidad de la aplicación es un aspecto
  clave del desarrollo de aplicaciones web debido al constante avance que sufre
  este tipo de tecnologías. Gracias a desarrollar la aplicación en base a un
  \cursiva{framework de desarrollo rápido de aplicaciones} conseguimos una alta
  mantenibilidad y legibilidad del \cursiva{código fuente} de la aplicación.

\item \negrita{Portabilidad}: La aplicación web, al estar desarrollada bajo
  \cursiva{Ruby On Rails} tiene una amplia portabilidad en cuánto al lado
  servidor ya que éste es compatible con los principales sistemas operativos del
  mercado.

  En cuanto al lado cliente tenemos la portabilidad asegurada ya que sólo será
  necesario que el usuario disponga de una conexión a \cursiva{internet} y un
  navegador web compatible con los estándares.
\end{itemize}

\section{Descripción de Requisitos Funcionales}

\subsection{Configurar el sitio}
El administrador podrá navegar por la página de configuración web y la
usará para cambiar el comportamiento de la misma.\\

Esta función se realizará raras veces.
\subsection{Crear una cuenta de usuario}
Para que un usuario cualquiera sea capaz de poder usar nuestro producto ha de
tener una cuenta registrada en el servidor que lo distinga unívocamente de
cualquier otro usuario; para ello se le pedirá introducir un nombre de usuario,
y una dirección de e-mail. Una vez registrado el sistema enviará un e-mail de
confirmación de
registro dando la bienvenida a nuestro sistema.\\

Esta función es condición necesaria para que puedan realizarse todos y cada uno
de los requisitos que se expresan en los apartados siguientes y se realizará una
única vez por usuario.

\subsection{Acceder a la cuenta de usuario}
Para que cualquier usuario registrado sea capaz de usar nuestro producto
deberemos de proporcionarle un método de acceso e identificación mediante el
nombre de cuenta y contraseña que han sido
albergados en nuestro servidor.\\

Esta operación se realizará muy a menudo y tantas veces como el usuario desee.
\subsection{Abandonar la cuenta de usuario}
Cualquier usuario que haya iniciado sesión en nuestra aplicación
deberá ser capaz de abandonar de forma segura la misma.\\

Esta operación se realizará muy a menudo y tantas veces como un usuario que haya
iniciado sesión desee.

\subsection{Abrir equipos}
El administrador será capaz de decidir cuándo los usuarios del sistema podrán
modificar sus alineaciones o tácticas, así como la posibilidad de interactuar
con los sistemas de finanzas y ofertas de jugadores.

Esta función será realizada muy a menudo.

\subsection{Cerrar equipos}
El administrador será capaz de decidir cuándo los usuarios del sistema no podrán
modificar sus alineaciones o tácticas, así como la posibilidad de interactuar
con los sistemas de finanzas y ofertas de jugadores.

Esta función será realizada muy a menudo.

\subsection{Listar ofertas}
El usuario podrá en cualquier momento visionar el listado de ofertas pendientes,
aceptadas o rechazadas que tenga en su historial.

Esta operación podrá ser realizada en cualquier momento.

\subsection{Emitir oferta por futbolista}
El usuario podrá en cualquier momento hacer una oferta por un futbolista en el
que esté interesado, interactuando de este modo con otros usuarios a través de
acuerdos económicos. Para evitar que dos usuarios puedan realizar acuerdos
económicos para hacer trampas y beneficiarse deportiva y económicamente, éstos
no podrán tasar a sus propios jugadores haciendo de la oferta una especie de
pujo, sino que será el propio sistema el que estipulará el precio del futbolista
en base a la calidad que éste disponga.

Esta operación podrá ser realizada en cualquier momento en el que los equipos
estén abiertos.

\subsection{Aceptar oferta por futbolista}
El usuario podrá en cualquier momento aceptar una oferta de un futbolista de su
propiedad emitida por otro usuario.

Esta operación podrá ser realizada en cualquier momento en el que los equipos
estén abiertos.

\subsection{Rechazar oferta por futbolista}
El usuario podrá en cualquier momento rechazar una oferta de un futbolista de su
propiedad emitida por otro usuario.

Esta operación podrá ser realizada en cualquier momento en el que los equipos
estén abiertos.

\subsection{Cancelar oferta por futbolista}
Un usuario que esté arrepentido de una oferta realizada por un futbolista podrá
cancelarla en cualquier momento.

Esta operación podrá ser realizada en cualquier momento en el que los equipos
estén abiertos.

\subsection{Renovación de futbolistas}
El sistema actualizará los sueldos que reciben los futbolistas al término de
cada temporada, según la ponderación de los atributos
cualitativos que tengan.\\

Esta operación se realizará una vez por temporada y afectará a todos los
futbolistas del sistema.
\subsection{Venta de entradas}
El usuario podrá estipular el precio de venta de las entradas en los partidos
que se disputen en su estadio siendo ésta la forma principal
de obtención de ingresos para la gestión económica del club.\\

Esta función será realizada tantas veces como desee el usuario.

\subsection{Planificación de alineaciones}
El usuario podrá escoger cuál será la alineación que presentará para el próximo
partido que vaya a disputar su equipo, pasando ésta a ser la alineación por
defecto hasta que el usuario decida volver a
cambiarla.\\

Esta operación se realizará a menudo.
\subsection{Planificación de tácticas}
El usuario podrá escoger entre diversas tácticas de juego para afrontar los
distintos partidos. Esta acción repercutirá directamente en el estilo de juego y
por tanto en la simulación resultado del
partido que dispute.\\

Esta operación se realizará tantas veces como el usuario desee.
\subsection{Planificación de entrenamientos}
El usuario podrá personalizar los atributos a entrenar de los futbolistas
integrantes de su equipo. De esta forma el club adquirirá
valor deportivo y económico.\\

Esta función se realizará a menudo.

\subsection{Asignar club}
Una vez que un usuario haya sido registrado en el servidor mediante una cuenta
de usuario el sistema le asignará un nuevo club cuando
todas las ligas hayan concluido y el administrador así desee.\\

Esta operación se realizará una vez por temporada y por cada nuevo usuario
registrado en el sistema.

\subsection{Simular jornadas}
El sistema debe emitir comentarios durante el transcurso del partido para que
cualquier usuario sea capaz de conocer la evolución del mismo.\\

Esta función se realizará a menudo, una vez por jornada.

\subsection{Comenzar jornadas}
El sistema debe simular la disputa de un partido en base a la ponderación de los
atributos cualitativos de los futbolistas alineados, tácticas empleadas y
variables aleatorias confrontadas con las del equipo
rival.\\

Esta función se realizará a menudo, una vez por jornada.

\subsection{Pasar a siguiente jornada}
El sistema debe ser capaz de gestionar el calendario creado de cada
competición. Para ello se le da la oportunidad al administrador de avanzar en la
línea temporal del calendario para hacer discurrir las jornadas del calendario.

Además, en esta acción se actualizarán los entrenamientos y traspasos aceptados
por futbolistas entre usuarios, así como restar los gastos de la plantilla y
mantenimiento del estadio.

Esta función se realizará a menudo a lo largo de la competición.

\subsection{Promoción y descenso de clubes}
El sistema debe ser capaz de gestionar varias ligas organizadas de forma
jerárquica según el número de usuarios que estén registrados en nuestro
servidor. De esta forma conseguiremos que el usuario desee mejorar la posición
de su club, lo que repercutirá directamente en la
capacidad deportiva y económica del mismo.\\

Esta función se realizará una vez por temporada y una vez por jornada.

\subsection{Composición de calendarios}
El sistema deberá de crear un calendario de partidos por cada liga que está
albergada en el servidor. Los calendarios se crearán una vez por
temporada y será accesible a todos los usuarios.\\

Esta función se realizará una vez por creación de liga.
\subsection{Gestiones de lesiones y sanciones}
El sistema debe ser capaz de llevar a cabo una política realista de imposición
de sanciones a los futbolistas, así como simular lesiones
más o menos duraderas.\\

Esta función se realizará a menudo.

\section{Especificación de los Casos de Uso}
\subsection{Caso de uso crear cuenta de usuario}
\negrita{Descripción: } Crea una cuenta de usuario para poder
participar y acceder a las acciones provistas por nuestra aplicación.\\
\negrita{Actores: } Usuario no registrado, usuario registrado,
sistema. \\
\negrita{Precondición: } Ninguna. \\
\negrita{Postcondición: } Se da de alta un nuevo usuario en el
sistema. \\
\negrita{Escenario principal:}
\begin{enumerate}
\item Un usuario no registrado desea registrarse.
\item El usuario introduce el nombre de usuario, la contraseña asociada al
  mismo, una cuenta de correo electrónico además de repetir la contraseña para
  evitar un error humano al ingresarla y también el idioma en el que desea ver
  por defecto la web y recibir los correos electrónicos.
\item El sistema registra al usuario.
\item El sistema avisa al usuario registrado mediante un correo electrónico y la
  propia interfaz web que ha completado su registro con éxito.
\end{enumerate}
\negrita{Escenario alternativo:}
\begin{enumerate}
\item Si el usuario introduce un nombre de cuenta de usuario o e-mail existente
  o dos contraseñas no coincidentes se le avisará de los errores cometidos para
  que pueda realizar un registro exitoso.
\end{enumerate}

\subsection{Caso de uso acceder a la cuenta de usuario}
\negrita{Descripción: } El usuario provee un sistema para poder identificar a
los usuarios registrados en el mismo y dejarlos acceder
a la funciones principales de nuestro sistema.\\
\negrita{Actores: } Usuario registrado, usuario logueado. \\
\negrita{Precondición: } El usuario ha de estar registrado en nuestro sistema. \\
\negrita{Postcondición: } Se identifica al usuario en el sistema y se
le permite acceder a las funciones principales del mismo. \\
\negrita{Escenario principal:}
\begin{enumerate}
\item Un usuario registrado desea iniciar sesión.
\item El usuario introduce el nombre de usuario y la contraseña asociada al
  mismo, además de repetir la contraseña para evitar un error humano al
  ingresarla.
\item El sistema registra al usuario.
\item El sistema avisa al usuario registrado mediante un correo electrónico y la
  propia interfaz web que ha completado su registro con éxito.
\end{enumerate}
\negrita{Escenario alternativo:}
\begin{enumerate}
\item Si el usuario introduce un nombre de cuenta de usuario o dos contraseñas
  no coincidentes se le avisará de los errores cometidos para que pueda realizar
  un registro exitoso.
\end{enumerate}

\subsection{Caso de uso abandonar la cuenta de usuario}
\negrita{Descripción: } El sistema provee un sistema para que un
usuario logueado pueda desconectarse de forma segura de nuestro sistema.\\
\negrita{Actores: } Usuario logueado. \\
\negrita{Precondición: } El usuario ha de haber iniciado sesión en
nuestro sistema. \\
\negrita{Postcondición: } Se cierra la sesión del usuario con nuestro sistema. \\
\negrita{Escenario principal:}
\begin{enumerate}
\item El usuario logueado desea abandonar su sesión.
\item Se le pregunta al usuario si está seguro de realizar la acción.
\item El sistema avisa al usuario que acaba de finalizar la sesión a través de
  la interfaz web de nuestra aplicación.
\end{enumerate}
\negrita{Escenario alternativo:}
\begin{enumerate}
\item Si el usuario decide que no está seguro de abandonar la sesión no se hace
  nada.
\end{enumerate}

\subsection{Caso de uso abrir equipos}
\negrita{Descripción: } El sistema provee un sistema para que el usuario
administrador pueda permitir cambios de jugadores, tácticas o
cualquier tipo de acción que modifique el estado de los equipos.\\
\negrita{Actores: } Usuario administrador. \\
\negrita{Precondición: } Ninguna. \\
\negrita{Postcondición: } Los equipos quedarán abiertos. \\
\negrita{Escenario principal:}
\begin{enumerate}
\item El administrador cierra los equipos.
\item El sistema avisa al administrador que acaba de abrir los equipos a través
  de la interfaz web de nuestra aplicación.
\end{enumerate}

\subsection{Caso de uso cerrar equipos}
\negrita{Descripción: } El sistema provee un sistema para que el usuario
administrador pueda evitar cambios de jugadores, tácticas o
cualquier tipo de acción que modifique el estado de los equipos.\\
\negrita{Actores: } Usuario administrador. \\
\negrita{Precondición: } Ninguna. \\
\negrita{Postcondición: } Los equipos quedarán cerrados. \\
\negrita{Escenario principal:}
\begin{enumerate}
\item El administrador cierra los equipos.
\item El sistema avisa al administrador que acaba de cerrar los equipos a través
  de la interfaz web de nuestra aplicación.
\end{enumerate}

\subsection{Caso de uso de notificación de registro por e-mail}
\negrita{Descripción: } El sistema notificará por correo electrónico
que un registro asociado al e-mail proporcionado se ha efectuado correctamente.\\
\negrita{Actores: } Sistema y usuario registrado \\
\negrita{Precondición: } El usuario ha de haber realizado un registro exitoso. \\
\negrita{Postcondición: } El sistema remite un correo electrónico en
el idioma que seleccionó en el registro nuestro usuario. \\
\negrita{Escenario principal:}
\begin{enumerate}
\item El sistema detecta que un nuevo usuario ha sido creado.
\item Remite un correo electrónico de bienvenida a la aplicación a la cuenta de
  correo asociada al usuario registrado.
\end{enumerate}

\subsection{Caso de uso de listar ofertas emitidas}
\negrita{Descripción: } El sistema ha de proporcionará una interfaz para poder
listar el historial de ofertas emitidas del club asociado.
\negrita{Actores: } Usuario logueado. \\
\negrita{Precondición: } El usuario ha de estar logueado en nuestra aplicación.
\negrita{Postcondición: } Se muestra un listado del historial de ofertas
emitidas del club.  \negrita{Escenario principal:}
\begin{enumerate}
\item Se muestra un listado del historial de ofertas emitidas del club.
\end{enumerate}

\subsection{Caso de uso de listar ofertas recibidas}
\negrita{Descripción: } El sistema ha de proporcionará una interfaz para poder
listar el historial de ofertas recibidas del club asociado.
\negrita{Actores: } Usuario logueado. \\
\negrita{Precondición: } El usuario ha de estar logueado en nuestra aplicación.
\negrita{Postcondición: } Se muestra un listado del historial de ofertas
recibidas del club.  \negrita{Escenario principal:}
\begin{enumerate}
\item Se muestra un listado del historial de ofertas recibidas del club.
\end{enumerate}

\subsection{Caso de uso de emitir oferta por futbolista}
\negrita{Descripción: } El sistema ha de proporcionará una interfaz para poder
realizar una oferta por un futbolista ajeno al club del que el usuario es
propietario.
\negrita{Actores: } Usuario logueado. \\
\negrita{Precondición: } El usuario ha de estar logueado en nuestra aplicación,
el comprador ha de tener la suficiencia económica
para acometer el fichaje, además los equipos han de estar abiertos. \\
\negrita{Postcondición: } El usuario emite una oferta por un futbolista ajeno a
su club y el propietario del mismo recibe la misma.  \negrita{Escenario
  principal:}
\begin{enumerate}
\item Un usuario emite una oferta por un futbolista que posee otro usuario.
\item El usuario propietario del futbolista recibe la oferta.
\end{enumerate}
\negrita{Escenario alternativo:}
\begin{enumerate}
\item El usuario comprador no puede efectuar la oferta.
\item El sistema emitirá un aviso a través de la interfaz web para informar que
  es imposible emitir la oferta.
\end{enumerate}

\subsection{Caso de uso de cancelar oferta emitida por futbolista}
\negrita{Descripción: } El sistema ha de proporcionará una interfaz para poder
cancelar una oferta emitida por un futbolista.
\negrita{Actores: } Usuario logueado. \\
\negrita{Precondición: } El usuario ha de estar logueado en nuestra aplicación,
la oferta ha de existir y los equipos han de estar
abiertos. \\
\negrita{Postcondición: } El usuario cancela una oferta por un futbolista de un
equipo ajeno al suyo.  \negrita{Escenario principal:}
\begin{enumerate}
\item El usuario cancela la oferta previamente realizada.
\end{enumerate}

\subsection{Caso de uso de rechazar oferta por futbolista}
\negrita{Descripción: } El sistema ha de proporcionará una interfaz para poder
rechazar una oferta realizada por un futbolista que provenga de otro usuario.
\negrita{Actores: } Usuario logueado. \\
\negrita{Precondición: } El usuario ha de estar logueado en nuestra aplicación,
la oferta ha de existir y los equipos han de estar
abiertos. \\
\negrita{Postcondición: } El usuario rechaza una oferta por un futbolista de un
equipo ajeno al suyo.  \negrita{Escenario principal:}
\begin{enumerate}
\item El usuario rechaza la oferta previamente realizada.
\end{enumerate}

\subsection{Caso de uso de aceptar oferta por futbolista}
\negrita{Descripción: } El sistema ha de proporcionará una interfaz para poder
aceptar una oferta realizada por un futbolista que provenga de otro usuario.
\negrita{Actores: } Usuario logueado. \\
\negrita{Precondición: } El usuario ha de estar logueado en nuestra aplicación,
la oferta ha de existir, los equipos han de estar abiertos, la venta del jugador
no debe suponer rebasar el límite
mínimo de jugadores en plantilla. \\
\negrita{Postcondición: } El usuario acepta la oferta, el montante económico de
la clausula del jugador es transferido a las arcas del club y el jugador ya no
es perteneciente a la plantilla.  \negrita{Escenario principal:}
\begin{enumerate}
\item El usuario acepta la oferta.
\item El futbolista es transferido a la plantilla del usuario emisor de la
  oferta.
\item Se incrementa la caja del club en el montante económico de la clausula del
  jugador transferido.
\end{enumerate}
\negrita{Primer Escenario alternativo: }
\begin{enumerate}
\item El usuario emisor no tiene recursos económicos suficientes para acometer
  el fichaje.
\item El sistema informará al usuario que no puede aceptar la oferta y ésta será
  automáticamente rechazada.
\end{enumerate}
\negrita{Segundo Escenario alternativo: }
\begin{enumerate}
\item El usuario receptor de la oferta está en el número mínimo de jugadores en
  plantilla.
\item El sistema informará al usuario que no puede aceptar la oferta y la misma
  será automáticamente rechazada.
\end{enumerate}
\subsection{Caso de uso de renovación de futbolistas}
\negrita{Descripción: } El sistema renovará cada final de temporada a los
jugadores de nuestro equipo actualizando el sueldo que reciben y
su clausula de rescisión.\\
\negrita{Actores: } Sistema.\\
\negrita{Precondición: } Las temporadas han de haber concluido y los
equipos bloqueados. \\
\negrita{Postcondición: } Los jugadores serán renovados en función de
la nueva ponderación de sus atributos así como su clausula de rescisión.\\
\negrita{Escenario principal:}
\begin{enumerate}
\item El sistema realizará una ponderación de los atributos de los futbolistas.
\item Se asigna una clausula de rescisión en función de la ponderación.
\item Se asigna un sueldo anual en función de la ponderación.
\end{enumerate}

\subsection{Caso de uso de fijar precio de venta de entradas}
\negrita{Descripción: } El usuario podrá modificar el precio de las entradas de
su equipo de fútbol. Este hecho repercutirá directamente en la afluencia al
estadio por parte de nuestros aficionados y a las
arcas del club.\\
\negrita{Actores: } Usuario logueado.\\
\negrita{Precondición: } Los equipos deben estar desbloqueados\\
\negrita{Postcondición: } El precio del ticket de entrada es
modificado.\\
\negrita{Escenario principal:}
\begin{enumerate}
\item El usuario desea cambiar el precio del ticket de entrada.
\item El ticket de entrada es cambiado al precio estipulado por el usuario.
\end{enumerate}
\negrita{Escenario alternativo:}
\begin{enumerate}
\item El usuario introduce un precio menor de cero.
\item El sistema informará al usuario que no puede asignar un precio negativo a
  través de la interfaz web.
\end{enumerate}

\subsection{Caso de uso de planificación de alineaciones}
\negrita{Descripción: } El usuario podrá modificar quiénes serán los
futbolistas que jueguen el siguiente partido.\\
\negrita{Actores: } Usuario logueado\\
\negrita{Precondición: } Los equipos deben estar desbloqueados.\\
\negrita{Postcondición: } El usuario dejará formada la alineación que
usará en el siguiente partido.\\
\negrita{Escenario principal:}
\begin{enumerate}
\item El usuario desea cambiar la posición entre dos jugadores.
\item El usuario seleccionará dos jugadores.
\item Efectuará el cambio de posición.
\item El usuario repetirá este proceso todas las veces que necesite hasta dejar
  la alineación como desee.
\end{enumerate}
\negrita{Escenario alternativo:}
\begin{enumerate}
\item El usuario selecciona un sólo jugador.
\item Efectúa el cambio de posición.
\item El sistema emitirá un aviso para informar que la operación que intenta
  hacer es imposible.
\end{enumerate}

\subsection{Caso de uso de planificación de tácticas}
\negrita{Descripción: } El usuario podrá seleccionar una táctica
distinta para afrontar el próximo partido.\\
\negrita{Actores: } Usuario logueado.\\
\negrita{Precondición: } Los equipos han de estar desbloqueados.\\
\negrita{Postcondición: } Se cambia la táctica por la seleccionada por
el usuario.\\
\negrita{Escenario principal:}
\begin{enumerate}
\item El usuario desea cambiar la táctica usada en el partido.
\item El usuario selecciona la táctica que desea usar y la fija como táctica por
  defecto para los siguientes partidos.
\end{enumerate}

\subsection{Caso de uso de planificación de entrenamientos}
\negrita{Descripción: } El usuario podrá mejorar una de las
habilidades que integran un jugador cada vez que lo mande a entrenar.\\
\negrita{Actores: } Usuario logueado.\\
\negrita{Precondición: } Los equipos han de estar desbloqueados y el
jugador no puede estar entrenando otra habilidad previa.\\
\negrita{Postcondición: } El jugador entrenará la habilidad
seleccionada por el usuario las jornadas que el sistema estipule.\\
\negrita{Escenario principal:}
\begin{enumerate}
\item El usuario desea que un jugador de su plantilla entrene cierta habilidad.
\item El sistema informará de cuántas jornadas serán necesarias para terminar el
  entrenamiento.
\end{enumerate}
\negrita{Escenario alternativo: }
\begin{enumerate}
\item Si el usuario intenta poner a entrenar la habilidad de un jugador que haya
  llegado al máximo posible para el atributo el sistema le informará a través de
  la interfaz web de que no es necesario realizar dicha acción.
\end{enumerate}

\subsection{Caso de uso asignar club}
\negrita{Descripción: } El sistema proporciona una forma de asignar nuevos club
de fútbol a los nuevos usuarios registrados cuando el administrador lo desee
entre
temporada y temporada.\\
\negrita{Actores: } Sistema, usuario registrado y administrador \\
\negrita{Precondición: } Las temporadas han de haber finalizado, debe haber un
mínimo de nuevos usuarios registrados para poder formar la
nueva liga y los equipos han de estar bloqueados. \\
\negrita{Postcondición: } Se les asigna a cada usuario un nuevo club con varios
jugadores y un monto inicial de dinero para poder comenzar
con las transacciones. \\
\negrita{Escenario principal:}
\begin{enumerate}
\item El administrador desea crear una nueva liga.
\item El sistema verifica que haya los suficientes nuevos usuarios sin clubes
  asignados para crearla.
\item El sistema genera una plantilla de futbolistas nuevos a cada nuevo club
  creado.
\item Se asigna cada club a cada usuario nuevo que formará la nueva liga.
\end{enumerate}

\subsection{Caso de uso de simular jornadas}
\negrita{Descripción: } El sistema simulará cada partido de cada jornada de cada
temporada según las alineaciones y tácticas de los
equipos que se enfrenten.\\
\negrita{Actores: } Administrador y Sistema. \\
\negrita{Precondición: } Los equipos deben estar bloqueados.\\
\negrita{Postcondición: } Se generará un conjunto de comentarios
ordenados por cada acción que tenga lugar en el partido.\\
\negrita{Escenario principal: }
\begin{enumerate}
\item El administrador decidirá simular los partidos de la jornada en la que se
  encuentre el sistema.
\item El sistema generará una ristra de comentarios por cada acción generada en
  el encuentro usando la ponderación de los equipos alineados y las tácticas
  utilizadas.
\end{enumerate}

\subsection{Caso de uso de comenzar jornadas}
\negrita{Descripción: } El sistema emulará la retransmisión de los
partidos englobados en las jornadas en tiempo real.\\
\negrita{Actores: } Administrador y Sistema. \\
\negrita{Precondición: } Los equipos deben estar bloqueados.\\
\negrita{Postcondición: } Se emitirá a través de la interfaz web el
conjunto de comentarios generados previamente en la simulación de los partidos.\\
\negrita{Escenario principal: }
\begin{enumerate}
\item El administrador decide que se comiencen a emitir los comentarios de los
  partidos.
\item El sistema empieza a mostrar los comentarios de cada partido de forma
  ordenada y actualizándose cada minuto.
\end{enumerate}

\subsection{Caso de uso de pasar a siguiente jornada}
\negrita{Descripción: } El sistema proveerá una forma de avanzar a lo
largo del calendario de partidos y actualizar traspasos, finanzas y entrenamientos.\\
\negrita{Actores: } Administrador y Sistema.\\
\negrita{Precondición: } Los equipos deben estar bloqueados\\
\negrita{Postcondición: }Se pasará de jornada en el calendario de las ligas, se
actualizarán los entrenamientos, traspasos y finanzas de los
clubes registrados en el sistema.\\
\negrita{Escenario principal: }
\begin{enumerate}
\item El administrador decide pasar a la siguiente jornada.
\item El sistema actualiza los resultados y las clasificaciones de cada liga y
  hace de la jornada siguiente la actual.
\item Se actualizan las finanzas de los clubes.
\item Se actualizan los traspasos aceptados.
\item Se actualizan los entrenamientos de los jugadores.
\end{enumerate}
\negrita{Escenario alternativo: }
\begin{enumerate}
\item Si la jornada es la última se dará por concluidas las temporadas que
  tenían lugar en cada liga.
\end{enumerate}


\subsection{Caso de uso de promoción y descenso de clubes}
\negrita{Descripción: } El sistema proveerá una forma de organizar las
ligas de forma jerárquica.\\
\negrita{Actores: } Administrador y Sistema.\\
\negrita{Precondición: } Los equipos deben estar bloqueados\\
\negrita{Postcondición: } Se ascenderán y descenderán los clubes que estén en
posiciones de descenso y ascenso según la organización
jerárquica de las ligas.\\
\negrita{Escenario principal: }
\begin{enumerate}
\item El sistema reorganizará las ligas para reflejar los cambios en las ligas
  debida a las posiciones que obtengan en la tabla de clasificación los equipos
  de nuestro sistema.
\end{enumerate}

\subsection{Caso de uso de composición de calendarios}
\negrita{Descripción: } El sistema proveerá una forma de componer los
calendarios de las ligas que se estén jugando en nuestra aplicación.\\
\negrita{Actores: } Administrador y Sistema.\\
\negrita{Precondición: } Los equipos deben estar bloqueados y las
temporadas finalizadas.\\
\negrita{Postcondición: } Se generará un calendario de jornadas siguiendo un
algoritmo round-robin de creación de calendarios por cada
liga albergada en nuestro sistema.\\
\negrita{Escenario principal: }
\begin{enumerate}
\item El sistema creará un calendario por liga de las nuevas temporadas que van
  a jugarse en nuestro sistema.
\end{enumerate}

\clearpage
\section{Diagramas de Casos de Uso}

\figura{casos_uso/funciones_principales.png}{scale=0.35}{Diagrama de casos de
  uso de las funciones principales}{casos-uso-funciones-generales}{H}

\negrita{Descripción:}
\begin{enumerate}
\item{Gestión de usuarios}
  \begin{itemize}
  \item Configurar el sitio
  \item Crear una cuenta de usuario
  \item Acceder a la cuenta de usuario
  \item Abandonar la cuenta de usuario
  \item Notificación de registro por e-mail
    \figura{casos_uso/gestion_usuarios.png}{scale=0.35}{Diagrama de casos de uso
      para la gestión de usuarios}{casos-uso-gestion-usuarios}{H}
  \end{itemize}
\item{Gestión económica}
  \begin{itemize}
  \item Listar ofertas emitidas
  \item Listar ofertas recibidas
  \item Emitir oferta por futbolista
  \item Aceptar oferta por futbolista
  \item Rechazar oferta por futbolista
  \item Cancelar oferta emitida por futbolista
  \item Renovación de futbolistas
  \item Fijar precio de venta de entradas
    \figura{casos_uso/gestion_economica.png}{scale=0.35}{Diagrama de casos de
      uso para la gestión económica}{casos-uso-gestion-economica}{H}
  \end{itemize}
\item{Gestión deportiva}
  \begin{itemize}
  \item Planificación de alineaciones
  \item Planificación de tácticas
  \item Planificación de entrenamientos
    \figura{casos_uso/gestion_deportiva.png}{scale=0.35}{Diagrama de casos de
      uso para la gestión deportiva}{casos-uso-gestion-deportiva}{H}
  \end{itemize}
\item{Gestión de campeonatos}
  \begin{itemize}
  \item Asignar club
  \item Simular jornadas
  \item Comenzar jornadas
  \item Pasar a siguiente jornada
  \item Promoción y descenso de clubes
  \item Composición de calendarios
    \figura{casos_uso/gestion_campeonatos.png}{scale=0.35}{Diagrama de casos de
      uso para la gestión de campeonatos}{casos-uso-gestion-campeonatos}{H}
  \end{itemize}
\end{enumerate}



\chapter{Diseño}
\label{cap:diseño}
% -*-cap5.tex-*- Este fichero es parte de la plantilla LaTeX para la realización
% de Proyectos Final de Carrera, protejido bajo los términos de la licencia
% GFDL.  Para más información, la licencia completa viene incluida en el fichero
% fdl-1.3.tex Copyright (C) 2009 Pablo Recio Quijano

Para la implementación de toda la funcionalidad descrita en el capítulo anterior
seguiremos las convenciones que el framework \cursiva{Ruby on Rails} nos
proporciona para una mayor velocidad y corrección en el desarrollo.

\section{Arquitectura}

\label{pattern:mvc}
Debido a \cursiva{Ruby on Rails} es un framework que hace un uso intensivo del
patrón \cursiva{Modelo-Vista-Controlador} (MVC) \cite{agilerails}.

El patrón \cursiva{Modelo-Vista-Controlador} divide el software en tres capas
bien diferenciadas que se adaptan muy bien a la tecnología
Cliente-Servidor. Dichas capas son:

\begin{itemize}
\item \negrita{Modelo}: Esta capa es la relativa a la representación de la
  información con la que el proyecto interactuará. Está compuesto por las clases
  necesarias para el acceso, modificación, búsqueda y demás operaciones con los
  datos.
\item \negrita{Vista}: En esta capa agruparemos todas las clases necesarias para
  la representación de la interfaz con la que nuestros usuarios interactuarán.
\item \negrita{Controlador}: Será la capa que nos permita gestionar las
  peticiones que lleguen a nuestro servidor.
\end{itemize}

\subsection{Capa Modelo}

Para la capa de modelo y la consecuente persistencia de datos vamos a usar el
\cursiva{ORM} (mapeo objeto-relacional) que nos proporciona por defecto
\cursiva{Ruby On Rails}, la clase \cursiva{ActiveRecord} que implementa el
patrón \cursiva{Active Record} pudiéndonos abstraer del \cursiva{SGBD}
relacional que usemos y las especificaciones únicas del lenguaje \cursiva{SQL}
que implemente el mismo.

\subsection{Capa Vista}

Para esta capa vamos a utilizar la clase \cursiva{ActionView}. Dicha clase es
una forma poderosa, sencilla y elegante, de embeber código \cursiva{Ruby} dentro
de cualquier tipo de fichero de la capa de presentación de nuestra página web
(\cursiva{JavaScript}, \cursiva{CSS}).

Además, haremos uso de los frameworks \cursiva{CoffeScript}
\cite{lang:coffescript} y \cursiva{SASS} \cite{lang:sass} para el desarrollo de
los ficheros \cursiva{JavaScript} y \cursiva{CSS} respectivamente.

\subsection{Capa Controlador}
En esta capa haremos uso de la clase \cursiva{ApplicationController}. Esta clase
nos ayudará a modelar de forma sencilla y fácil el enrutador de nuestra
aplicación y las respuestas que se hagan a través del cliente en la petición
web.

Además, nos ayudará a establecer un flujo sencillo entre nuestras capas Vista y
Modelo consiguiendo un menor acoplamiento y una mayor cohesión de nuestro
sistema.

\subsection{Active Record}
Se basa en el patrón de mismo nombre el cual es un patrón de arquitectura que
podemos encontrar en software que almacene información en bases de datos
relacionales. Conforme a este patrón, un objeto que lo use habrá de integrar
funciones para insertar, actualizar, eliminar y otras que correspondan en mayor
o menor medida a las operaciones con columnas de la tabla de una base de datos.

Este patrón crea una interfaz que envuelve una tabla de una base de datos o una
vista en una clase. En base a ello, un objeto de la clase corresponde a una fila
en la tabla. La clase que envuelve a la tabla o vista implementa métodos o
propiedades para acceder a cada columna.

Este patrón es normalmente usado para dar persistencia a los datos en un
proyecto software, haciendo de \cursiva{ORM} para la aplicación. Normalmente,
las relaciones de clave foránea entre clases son implementadas como una
propiedad del objeto \cite{pattern:activerecord}.

Las funcionalidades que implementa la clase \cursiva{Active Record} en
\cursiva{Rails} son:

\begin{itemize}
\item \negrita{Creación}: Permite la creación de nuevas líneas en la base de
  datos.
\item \negrita{Condiciones}: Permite consultar datos en la base de datos usando
  para ello cadenas de texto o \cursiva{hashes} que representen la parte
  \cursiva{WHERE} en una sentencia SQL.
\item \negrita{Sobrescribir métodos de accesibilidad por defecto}: Se pueden
  sobrescribir los métodos de accesibilidad para modificar el comportamiento en
  las operaciones de escritura, creación, lectura, y demás operaciones de
  lectura/escritura.
\item \negrita{Métodos de consulta como atributos}: Se pueden crear nuevos
  métodos que se comporten como atributos para conocer si un atributo dentro de
  la tabla está inicializado.
\item \negrita{Posibilidad de acceder a los atributos antes de que hayan pasado
    el casting de tipo}: Es una funcionalidad orientada sobre todo en las
  situaciones de validación.
\item \negrita{Buscadores dinámicos en base al atributo}: Métodos estándar
  creados en base al nombre de los atributos correspondientes al nombre de la
  columna dentro de la tabla.
\item \negrita{Serializar vectores, hashes y otros objetos en columnas de
    texto}: Es posible guardar este tipo de objetos en formato texto usando el
  lenguaje \cursiva{YAML}.
\item \negrita{Herencia simple de tabla}: Es posible el uso de la herencia
  simple entre tablas.
\item \negrita{Conexión a múltiples bases de datos en diferentes modelos}: Es
  posible establecer diferentes conexiones a diferentes \cursiva{SGBD} para
  diferentes clases.
\end{itemize}

\subsection{Action View}
\cursiva{Action View} hace uso del sistema de plantillas \cursiva{eRuby} para
embeber código \cursiva{Ruby} dentro de un documento de texto. Es ampliamente
utilizado para dentro de documentos \cursiva{HTML} y \cursiva{JavaScript}.

\subsection{Action Controller}
Es el núcleo de una petición web en \cursiva{Ruby on Rails}. Está formado por
una o más acciones que son ejecutadas cuando se produzca una petición, justo
después se puede redireccionar a otra acción o renderizar una plantilla. Una
acción se define como un método público dentro del controlador, el cuál es
automáticamente accesible al servidor web a través del enrutador de
\cursiva{Rails}.

Por defecto, sólo la clase \cursiva{ApplicationController} dentro de una
aplicación \cursiva{Rails} hereda de la clase
\cursiva{ActionController:Base}. Los demás controladores heredarán de la clase
\cursiva{ApplicationController}. Este tipo de estructura nos proporciona una
clase para configurar varias cosas que sean comunes a cualquier petición web.

Esta clase nos proporciona además las siguientes funcionalidades:

\begin{itemize}
\item \negrita{Peticiones}: El objeto de petición es totalmente accesible y se
  usa principalmente para conocer la cabecera de la petición \cursiva{HTTP}.
\item \negrita{Parámetros}: Todos los parámetros que lleva la petición, ya sean
  desde \cursiva{GET} o \cursiva{POST} o de la \cursiva{URL}, son accesibles a
  través del \cursiva{método} \negrita{params} que devolverá un \cursiva{hash}
  uni o multidimensional con todos los parámetros asociados a la petición.
\item \negrita{Sesiones}: Las sesiones permiten almacenar objetos entre petición
  y petición. Son bastante útiles para objetos que no estén aún preparados para
  pasar a ser almacenados en la base de datos. Para acceder a las sesiones
  bastará con llamar al método \negrita{session} y este devolverá un
  \cursiva{hash} uni o multidimensional con las sesiones almacenadas.
\item \negrita{Respuestas}: Toda acción lleva asociada una respuesta, la cual
  alberga las cabeceras y el documento para ser mandado al navegador del
  usuario. El objeto respuesta es generado automáticamente a través del uso de
  \cursiva{renders} y redirecciones, con lo que no necesita la interacción del
  usuario.
\item \negrita{Renderizaciones}: La clase \cursiva{ActionController} manda
  contenido al usuario a través del uso de métodos de renderización. El más
  común es el uso de una plantilla para la renderización del contenido, a través
  de plantillas \cursiva{ERB}.
\item \negrita{Redirecciones}: Son usadas para moverse de una acción a otra.
\item \negrita{Redirecciones y renderizaciones múltiples}: Es posible renderizar
  y redireccionar varias veces en la misma petición.
\end{itemize}

\subsection{Diagrama de componentes y estructura}

Podemos ver el diagrama de componentes en la figura \ref{diag:componentes}.

\figura{mvc-rails.png}{scale=0.50}{Diagrama de componentes}{diag:componentes}{H}

El framework \cursiva{Ruby on Rails} hace especial hincapié en tres conceptos
muy básicos que caracterizan toda su estructura, que son:

\begin{itemize}
\item \negrita{Convention over Configuration (COC)}: Convención frente a
  configuración.
\item \negrita{Don't Repeat Yourself (DRY)}: No te repitas
\item \negrita{Keep It Simple (KIS)}: Mantenlo simple.
\end{itemize}

Gracias a ello seguiremos la estructura de directorios en árbol para nuestro
código fuente que genera \cursiva{Rails} para una aplicación vacía.

\begin{itemize}
\item \texttt{/pfc\_sfo/} raíz del proyecto \cursiva{Rails}.
\item \texttt{/pfc\_sfo/app/} ficheros fuente principales.
\item \texttt{/pfc\_sfo/app/controllers/} ficheros fuente para los
  controladores.
\item \texttt{/pfc\_sfo/app/helpers/} ficheros fuente para los métodos de ayuda
  para refactorización de códigos en las vistas.
\item \texttt{/pfc\_sfo/app/mailers/} ficheros fuente para el código fuente de
  los correos que se envían desde el servidor.
\item \texttt{/pfc\_sfo/app/models/} ficheros fuente correspondientes a los
  modelos de datos.
\item \texttt{/pfc\_sfo/app/views/} ficheros fuente para las vistas de la
  aplicación, incluyendo como parte de las mismas las de los correos.
\item \texttt{/pfc\_sfo/app/assets/} ficheros fuente que han de ejecutarse en
  el navegador, archivos \cursiva{CSS}, \cursiva{JavasScript}, imágenes, etc...
\item \texttt{/pfc\_sfo/config/}, ficheros fuente de configuración de la
  aplicación web.
\item \texttt{/pfc\_sfo/config/environments/} ficheros fuente de configuración
  de los distintos entornos, desarrollo, test y producción.
\item \texttt{/pfc\_sfo/config/initializers/} ficheros de inicialización del
  servidor.
\item \texttt{/pfc\_sfo/config/locales/} ficheros \cursiva{YAML} para los
  textos de traducción en diferentes idiomas.
\item \texttt{/pfc\_sfo/db/} ficheros autogenerados dónde se albergan la
  configuración de la Base de Datos.
\item \texttt{/pfc\_sfo/db/migrate/} ficheros para la manipulación de las
  tablas de la base de datos.
\item \texttt{/pfc\_sfo/lib/} ficheros para la extensión de la aplicación con
  otros subsistemas software.
\item \texttt{/pfc\_sfo/log/} ficheros log del servidor.
\item \texttt{/pfc\_sfo/public/} ficheros de acceso público a través del
  enrutador.
\item \texttt{/pfc\_sfo/script/} scripts para la generación automática de
  código y la automatización de procesos.
\item \texttt{/pfc\_sfo/test/} ficheros para almacenar test automáticos.
\end{itemize}

\subsection{Entorno de ejecución}
En la figura \ref{fig:entornoejecucion} podemos ver el diagrama de despliegue de
la aplicación puesta en marcha en cualquiera de los entornos, ya sea este de
desarrollo, test o producción.

El sistema (servidor) interactuará con uno o más clientes web (navegadores). El
servidor, a su vez estará formado por un servidor de aplicaciones web, en
nuestro caso \cursiva{Webrick}, y el sistema de gestión de base de datos será
\cursiva{SQLite3}.

Gracias a esta sencilla configuración no tendremos que configurar ninguna cuenta
ni servicio en nuestro servidor web.

\clearpage
\section{Diagramas de clases conceptuales}
En la figura \ref{diag:clases_modelo} podemos observar el diagrama de clases de
la capa modelo de nuestra aplicación según la recolección de requisitos que
hemos obtenido en el Capítulo \ref{cap:analisis} de esta memoria. Dichos
diagramas serán realizados en notación \cursiva{UML} \cite{uml:distilled_standard}.

\figura{models_complete.png}{scale=0.25}{Diagrama de Clases UML de la capa
  modelo}{diag:clases_modelo}{H}

\clearpage

En la figura \ref{diag:clases_controlador} podemos observar el diagramas de
clases de la capa controlador de nuestra aplicación según la recolección de
requisitos que hemos obtenido en el Capítulo \ref{cap:analisis}.

\figura{controllers_complete.png}{scale=0.22,angle=270}{Diagrama de Clases UML
  de la capa controlador}{diag:clases_controlador}{H}

\section{Diagramas de secuencia de sistemas}

En las siguientes secciones podremos observar los diagramas de secuencias del
sistema software que hemos desarrollado.

\subsection*{Diagrama de secuencia de crear cuenta de usuario}

En la figura \ref{diagrama-secuencia-crear-cuenta-usuario} podemos observar como
se construye un nuevo usuario dentro del sistema con los datos proporcionados
por el usuario web.

\figura{diagramas_secuencia/crear_cuenta_usuario.png}{scale=0.42}{Diagrama de
  secuencia de crear cuenta de
  usuario}{diagrama-secuencia-crear-cuenta-usuario}{H}
\newpage

\subsection*{Diagrama de secuencia de acceder a cuenta de usuario}

En la figura \ref{diagrama-secuencia-acceder-cuenta-usuario} podemos observar
como se consigue acreditar las credenciales (usuario y contraseña) para iniciar
sesión.

\figura{diagramas_secuencia/acceder_cuenta_usuario.png}{scale=0.49}{Diagrama de
  secuencia de acceder a cuenta de
  usuario}{diagrama-secuencia-acceder-cuenta-usuario}{H}
\newpage

\subsection*{Diagrama de secuencia de abandonar cuenta de usuario}

En la figura \ref{diagrama-secuencia-abandonar-cuenta-usuario} podemos observar
cómo se destruye la sesión iniciada por un usuario.

\figura{diagramas_secuencia/abandona_cuenta_usuario.png}{scale=0.59}{Diagrama de
  secuencia de abandonar cuenta de
  usuario}{diagrama-secuencia-abandonar-cuenta-usuario}{H}
\newpage

\subsection*{Diagrama de secuencia de abrir equipos}

En la figura \ref{diagrama-secuencia-abrir-equipos} cómo se establece el estado
de las temporadas para permitir operaciones por parte de los usuarios en los
equipos que poseen.

\figura{diagramas_secuencia/abrir_equipos.png}{scale=0.78}{Diagrama de secuencia
  de abrir equipos}{diagrama-secuencia-abrir-equipos}{H}

\subsection*{Diagrama de secuencia de cerrar equipos}
En la figura \ref{diagrama-secuencia-abrir-equipos} podemos observar cómo se
establece el estado de las temporadas para impedir que los usuarios puedan
realizar operaciones con sus equipos.

Este diagrama es igual al \ref{diagrama-secuencia-abrir-equipos} pero con la
ruta \cursiva{/leagues/close\_teams} y la llamada al modelo con el método
\cursiva{close}.


\subsection*{Diagrama de secuencia de notificación de registro por e-mail}
En la figura \ref{diagrama-secuencia-notificacion-registro} se observa cómo se
envía un correo al usuario una vez se ha registrado en la aplicación web.

\figura{diagramas_secuencia/crear_cuenta_usuario.png}{scale=0.43}{Diagrama de
  secuencia de notificación de registro por
  e-mail}{diagrama-secuencia-notificacion-registro}{H}


\subsection*{Diagrama de secuencia de listar ofertas emitidas}
En la figura \ref{diagrama-secuencia-listar-ofertas-emitidas} observamos cómo se
establece el filtro de búsqueda para mostrar el listado de ofertas emitidas.

\figura{diagramas_secuencia/listar_ofertas_emitidas.png}{scale=0.52}{Diagrama de
  secuencia de listar ofertas
  emitidas}{diagrama-secuencia-listar-ofertas-emitidas}{H}

\subsection*{Diagrama de secuencia de listar ofertas recibidas}
En la figura \ref{diagrama-secuencia-listar-ofertas-emitidas} observamos cómo se
establece el filtro de búsqueda para mostrar el listado de ofertas recibidas.

Este diagrama es igual al \ref{diagrama-secuencia-listar-ofertas-emitidas} pero
con la ruta \cursiva{/clubs/*/received\_offers} y la llamada al modelo club con
el método \cursiva{offers\_as\_seller}.

\subsection*{Diagrama de secuencia de emitir oferta por futbolista}
En la figura \ref{diagrama-secuencia-emitir-oferta-futbolista} se observa cómo
se crea una nueva instancia de oferta y se relaciona con los usuarios y el
futbolista involucrado en la misma.

\figura{diagramas_secuencia/emitir_oferta_futbolista.png}{scale=0.45}{Diagrama
  de secuencia de emitir oferta por
  futbolista}{diagrama-secuencia-emitir-oferta-futbolista}{H}

\subsection*{Diagrama de secuencia de cancelar oferta emitida por futbolista}
En la figura \ref{diagrama-secuencia-cancelar-oferta-emitida-futbolista} se
observa cómo se recupera la instancia de la oferta emitida y se modifica su
estado para cancelarla.

\figura{diagramas_secuencia/cancelar_oferta_emitida_futbolista.png}{scale=0.68}{Diagrama
  de secuencia de cancelar oferta emitida por
  futbolista}{diagrama-secuencia-cancelar-oferta-emitida-futbolista}{H}

\subsection*{Diagrama de secuencia de rechazar oferta recibida por futbolista}
En la figura \ref{diagrama-secuencia-rechazar-oferta-recibida-futbolista} se
observa cómo se recupera la instancia de la oferta recibida y se modifica su
estado para rechazar la misma.

\figura{diagramas_secuencia/rechazar_oferta_recibida_futbolista.png}{scale=0.68}{Diagrama
  de secuencia de rechazar oferta recibida por
  futbolista}{diagrama-secuencia-rechazar-oferta-recibida-futbolista}{H}

\newpage

\subsection*{Diagrama de secuencia de aceptar oferta recibida por futbolista}
En la figura \ref{diagrama-secuencia-aceptar-oferta-recibida-futbolista} se
observa cómo se recupera la instancia de la oferta recibida y se actualizan los
modelos de datos involucrados (futbolista, club y usuarios) para llevar a cabo
la transferencia del futbolista entre los clubes.

\figura{diagramas_secuencia/aceptar_oferta_recibida_futbolista.png}{scale=0.65}{Diagrama
  de secuencia de aceptar oferta recibida por
  futbolista}{diagrama-secuencia-aceptar-oferta-recibida-futbolista}{H}
\newpage

\subsection*{Diagrama de secuencia de renovación de futbolistas}
En la figura \ref{diagrama-secuencia-renovacion-futbolistas} se describe cómo se
renuevan automáticamente los futbolistas para las siguientes temporadas
actualizando sus atributos.

\figura{diagramas_secuencia/renovacion_futbolistas.png}{scale=0.63}{Diagrama de
  secuencia de renovación de
  futbolistas}{diagrama-secuencia-renovacion-futbolistas}{H}

El proceso de renovación está integrado en start.
\newpage

\subsection*{Diagrama de secuencia de fijar precio de venta de entradas}
En la figura \ref{diagrama-secuencia-fijar-precio-venta-entradas} se describe
cómo se modifica el atributo del precio de venta de entradas.

\figura{diagramas_secuencia/fijar_precio_venta_entradas.png}{scale=0.54}{Diagrama
  de secuencia de fijar precio de venta de
  entradas}{diagrama-secuencia-fijar-precio-venta-entradas}{H}

\subsection*{Diagrama de secuencia de planificación de alineaciones}
En la figura \ref{diagrama-secuencia-planificacion-alineaciones} se describe
cómo se cambian las alineaciones para el próximo partido a través de la
modificación del atributo de posición.

\figura{diagramas_secuencia/planificacion_alineaciones.png}{scale=0.52}{Diagrama
  de secuencia de emitir planificación de
  alineaciones}{diagrama-secuencia-planificacion-alineaciones}{H}

\subsection*{Diagrama de secuencia de planificación de tácticas}
En la figura \ref{diagrama-secuencia-planificacion-tacticas} se muestra cómo se
cambia la táctica usada para el próximo partido a través de la modificación del
atributo de táctica del club.

\figura{diagramas_secuencia/planificacion_tacticas.png}{scale=0.70}{Diagrama de
  secuencia de planificación de tácticas
}{diagrama-secuencia-planificacion-tacticas}{H}

\subsection*{Diagrama de secuencia de planificación de entrenamientos}
En la figura \ref{diagrama-secuencia-planificacion-entrenamientos} se muestra
cómo se crean nuevas instancias para la planificación de los entrenamientos de
los futbolistas asociados a un club.

\figura{diagramas_secuencia/planificacion_entrenamientos.png}{scale=0.45}{Diagrama
  de secuencia de planificación de entrenamientos
}{diagrama-secuencia-planificacion-entrenamientos}{H}

\subsection*{Diagrama de secuencia de asignar club}
En la figura \ref{diagrama-secuencia-asignar-club} se observa cómo se crea una
nueva instancia de un club para todo usuario que no tenga club asignado.

\figura{diagramas_secuencia/asignar_club.png}{scale=0.45}{Diagrama de secuencia
  de secuencia de asignar club }{diagrama-secuencia-asignar-club}{H}

\subsection*{Diagrama de secuencia de simular jornadas}
En la figura \ref{diagrama-secuencia-simular-jornadas} se observa cómo se crean
nuevas instancias de detalles de partidos por cada uno de los partidos no
jugados y pertenecientes a la jornada actual del sistema.

\figura{diagramas_secuencia/simular_jornadas.png}{scale=0.35,
  angle=270}{Diagrama de secuencia de simular
  jornadas}{diagrama-secuencia-simular-jornadas}{H}

\subsection*{Diagrama de secuencia de comenzar jornadas}
En la figura \ref{diagrama-secuencia-comenzar-jornadas} se observa cómo se
actualiza el atributo correspondiente de todos los partidos actuales para
estipular que el partido ha comenzado a la fecha actual del sistema.
\newpage

\figura{diagramas_secuencia/comenzar_jornadas.png}{scale=0.49,
  angle=270}{Diagrama de secuencia de comenzar
  jornadas}{diagrama-secuencia-comenzar-jornadas}{H}


\subsection*{Diagrama de secuencia de pasar a siguiente jornada}
En la figura \ref{diagrama-secuencia-pasar-siguiente-jornada} se detalla cómo se
actualizan los modelos de datos de temporada, jornada, jugador, entrenamiento,
finanzas y club para avanzar en la competición.
\newpage
 
\figura{diagramas_secuencia/pasar_siguiente_jornada.png}{scale=0.29,
  angle=270}{Diagrama de secuencia de pasar a siguiente
  jornada}{diagrama-secuencia-pasar-siguiente-jornada}{H}

\subsection*{Diagrama de secuencia de promoción y descenso de clubes}
En la figura \ref{diagrama-secuencia-promocion-descenso-clubes} se observa cómo
se crean nuevas instancias de nuevas temporadas y se modifican los estados de
las temporadas actuales para simular el avance en la competición y el progreso y
descenso jerárquico en las divisiones del juego.

\figura{diagramas_secuencia/promocion_descenso_clubes.png}{scale=0.38}{Diagrama
  de secuencia de promocion y descenso de
  clubes}{diagrama-secuencia-promocion-descenso-clubes}{H}


\subsection*{Diagrama de secuencia de composición de calendarios}
En la figura \ref{diagrama-secuencia-composicion-calendarios} se observa cómo se
crean nuevas instancias de las jornadas que estarán relacionadas con los clubes
en los que se coincida en las mismas divisiones.
\newpage

\figura{diagramas_secuencia/crear_calendario.png}{scale=0.36,
  angle=270}{Diagrama de secuencia de composición de
  calendarios}{diagrama-secuencia-composicion-calendarios}{H}


\chapter{Implementación}
\label{cap:implementacion}
% -*-cap6.tex-*- Este fichero es parte de la plantilla LaTeX para la realización
% de Proyectos Final de Carrera, protejido bajo los términos de la licencia
% GFDL.  Para más información, la licencia completa viene incluida en el fichero
% fdl-1.3.tex Copyright (C) 2009 Pablo Recio Quijano

En este capítulo vamos a describir y explicar los detalles de implementación del
sistema software; además de comentar los problemas y contingencias surgidas
durante el proceso de desarrollo del mismo.

Para situar mejor al lector pasaremos a descargar e instalar el código fuente
base de la aplicación desarrollada que será descargable desde
\cite{cod:pfc_sfo}.

El código fuente está albergado en un repositorio de tipo \cursiva{GIT}, aunque
también será descargable como un \cursiva{tarball} a través de la interfaz web.

Para que la aplicación sea totalmente funcional habrá que añadir los siguientes
directorios al árbol base de la aplicación:

\begin{itemize}
\item \texttt{/tmp/}
\item \texttt{/log/}
\end{itemize}

Para que el sistema pueda ser ejecutado, además se habrán de cumplir los
requisitos especificados en \ref{sec:restricciones_generales}. También será
importante tener configurado un SGBD relacional a escoger entre:

\begin{itemize}
\item MySQL
\item Postgres
\item SQLite3
\item Oracle
\end{itemize}

Una vez que tengamos el código fuente aislado en un directorio habremos de
ejecutar los siguientes comandos en orden:

\begin{enumerate}
\item \texttt{gem install}
\item \texttt{rake db:create}
\item \texttt{rake db:migrate}
\item \texttt{rails s}
\end{enumerate}

Con todos estos pasos dados ya podremos acceder a nuestra aplicación web a
través del navegador y de la dirección \url{http://localhost:3000/}.

\section{Modelo de datos}
Al usar como gestor \cursiva{ORM}, la clase incorporada en \cursiva{Rails},
\cursiva{ActiveRecord}, habrá que especificar los denominados ficheros de
migración para ir creando las tablas a través de código ruby. Vemos aquí un
ejemplo de uno de los archivos de migración extraídos de nuestro proyecto de fin
de carrera:

\lstinputlisting[style=Ruby]{codigo/20090401160020_create_players.rb}

Como podemos observar, a través de los bloques de Ruby creamos comandos para
inicializar y revertir los cambios que traen la migración consigo. Los cambios
que podemos realizar entre otras son:

\begin{itemize}
\item Crear tablas.
\item Eliminar tablas.
\item Crear columnas.
\item Eliminar columnas.
\item Modificar nombres de tablas.
\item Modificar nombres de columnas.
\item Insertar nuevos datos.
\item Eliminar datos.
\item Migrar datos.
\end{itemize}

El conjunto de ficheros de migración crearán el archivo \texttt{schema.rb} que
tiene una definición completa de las tablas que albergan nuestra base de datos:

\lstinputlisting[style=Ruby]{codigo/schema.rb}

\section{Controlador de versiones}
Para controlar los diversos cambios que realizamos sobre el producto o la
configuración del mismo hemos decidido usar \cursiva{GIT} el cuál se integra a
la perfección con nuestro entorno de desarrollo \cursiva{Linux} y además es uno
de los controladores de la versión más en boga del actual panorama software.

El sistema de versiones \cursiva{GIT} nos proporciona las siguientes
funcionalidades \cite{wiki:GIT}:

\begin{itemize}
\item Fuerte apoyo al desarrollo no-lineal, por ende rapidez en la gestión de
  ramas y mezclado de diferentes versiones. \cursiva{Git} incluye herramientas
  específicas para navegar y visualizar un historial de desarrollo
  no-lineal. Una presunción fundamental en \cursiva{Git} es que un cambio será
  fusionado mucho más frecuentemente de lo que se escribe originalmente,
  conforme se pasa entre varios programadores que lo revisan.
\item Gestión distribuida. Al igual que \cursiva{Darcs}, \cursiva{BitKeeper},
  \cursiva{Mercurial}, \cursiva{SVK}, \cursiva{Bazaar} y \cursiva{Monotone},
  \cursiva{Git} le da a cada programador una copia local del historial del
  desarrollo entero, y los cambios se propagan entre los repositorios
  locales. Los cambios se importan como ramas adicionales y pueden ser
  fusionados en la misma manera que se hace con la rama local.
\item Los almacenes de información pueden publicarse por \cursiva{HTTP},
  \cursiva{FTP}, \cursiva{rsync} o mediante un protocolo nativo, ya sea a través
  de una conexión \cursiva{TCP/IP} simple o a través de cifrado
  \cursiva{SSH}. \cursiva{Git} también puede emular servidores \cursiva{CVS}, lo
  que habilita el uso de clientes \cursiva{CVS} pre-existentes y modulos
  \cursiva{IDE} para \cursiva{CVS} pre-existentes en el acceso de repositorios
  \cursiva{Git}.
\item Los repositorios \cursiva{Subversion} y \cursiva{svk} se pueden usar
  directamente con \cursiva{git-svn}.
\item Gestión eficiente de proyectos grandes, dada la rapidez de gestión de
  diferencias entre archivos, entre otras mejoras de optimización de velocidad
  de ejecución.
\item Todas las versiones previas a un cambio determinado, implican la
  notificación de un cambio posterior en cualquiera de ellas a ese cambio
  (denominado autenticación criptográfica de historial). Esto existía en
  \cursiva{Monotone}.
\item Resulta algo más caro trabajar con ficheros concretos frente a proyectos,
  eso diferencia el trabajo frente a \cursiva{CVS}, que trabaja con base en
  cambios de fichero, pero mejora el trabajo con afectaciones de código que
  concurren en operaciones similares en varios archivos.
\item Los renombrados se trabajan basándose en similitudes entre ficheros,
  aparte de nombres de ficheros, pero no se hacen marcas explícitas de cambios
  de nombre con base en supuestos nombres únicos de nodos de sistema de
  ficheros, lo que evita posibles, y posiblemente desastrosas, coincidencias de
  ficheros diferentes en un único nombre.
\item Realmacenamiento periódico en paquetes (ficheros). Esto es relativamente
  eficiente para escritura de cambios y relativamente ineficiente para lectura
  si el reempaquetado (con base en diferencias) no ocurre cada cierto tiempo.
\end{itemize}

A continuación nos dispondremos a visualizar algunas estadísticas de nuestro
sistema software a través de la aplicación \cursiva{GitStats}
\cite{prog:gitstats} y el análisis de estadísticas de código incluido dentro de
\cursiva{Ruby On Rails}.

En la figura \ref{fig:estadisticas_codigo} podemos ver las estadísticas de
código clasificadas según extensión y ordenados por el número de líneas que
conforman la totalidad del proyecto.

\figura{gitstats.png}{scale=0.5}{Estadísticas del
  repositorio}{fig:estadisticas_codigo}{H}

En la figura \ref{fig:estadisticas_rails} podemos ver un análisis del código
fuente según la estructura de un proyecto realizado en \cursiva{Ruby On Rails}.

\figura{stats.png}{scale=0.5}{Estadísticas del código fuente
  Rails}{fig:estadisticas_rails}{H}

A continuación pasamos a describir los términos que aparecen en la figura
\ref{fig:estadisticas_rails} para una mejor interpretación de la misma:

\begin{itemize}
\item Lines: Son las líneas de código totales de cada una de las categorías.
\item LOC: Son las líneas de código reales no autogeneradas por el framework.
\item Classes: Son el número de clases pertenecientes a cada una de las
  categorías.
\item Methods: El número de métodos correspondientes a cada una de las
  categorías.
\item M/C: Media de métodos por clase.
\item LOC/M: Media de número de líneas por método.
\end{itemize}

\section{Herramientas de desarrollo}
\subsection{Codificación}
Para la creación y edición de los ficheros fuente del proyecto hemos utilizado
el editor de texto gratuito \cursiva{SublimeText2} \cite{prog:sublimetext2}, el
cual es lo bastante extensible para ser utilizado como un completo \cursiva{IDE}
de desarrollo para \cursiva{Ruby On Rails}.

Haremos intensivo uso de sus características avanzadas de editor de texto y de
autocomplección de código, así como de la característica de inserción rápida de
texto a través de los llamados \cursiva{snippets} \cite{prog:snippets}.

Otra característica avanzada de este editor de texto extensible es su
integración con \cursiva{GIT} y cualquier tecnología a través del uso de
plugins. Además de ser multiplataforma lo cual es idóneo para la depuración de
código \cursiva{CSS} en distintas plataformas.

\subsection{Depuración}
Para la depuración del código fuente usamos de forma extensiva el navegador web
\cursiva{Firefox} \cite{prog:firefox} el cuál nos proporciona la extensión
\cursiva{Firebug} \cite{prog:firebug}.

\cursiva{Firebug} es una extensión de cursiva{Firefox} creada y diseñada
especialmente para desarrolladores y programadores web. Es un paquete de
utilidades con el que se puede analizar (revisar velocidad de carga, estructura
\cursiva{DOM}), editar, monitorizar y depurar el código fuente, \cursiva{CSS},
\cursiva{HTML} y \cursiva{JavaScript} de una página web de manera instantánea e
inline.


\cursiva{Firebug} no es un simple inspector como \cursiva{Dom Inspector}, además edita y
permite guardar los cambios, un paso por delante del conocido \cursiva{Web
  Developer}. Su atractiva e intuitiva interfaz, con solapas específicas para el
análisis de cada tipo de elemento (consola, \cursiva{HTML}, \cursiva{CSS},
\cursiva{Script}, \cursiva{DOM} y red), permite al usuario un manejo fácil y
rápido. \cursiva{Firebug} está encapsulado en forma de plug-in o complemento de
\cursiva{Mozilla}, es \cursiva{Open Source}, libre y de distribución gratuita \cite{wiki:firebug}.


\backmatter % Apéndices, bibliografia ...


\clearpage
\addcontentsline{toc}{chapter}{Bibliografia y referencias}
\bibliographystyle{plain}
\bibliography{bibliografia}


\addcontentsline{toc}{chapter}{Software usado}
\chapter*{Software utilizado}
% -*-programas.tex-*-
% Este fichero es parte de la plantilla LaTeX para
% la realización de Proyectos Final de Carrera, protejido
% bajo los términos de la licencia GFDL.
% Para más información, la licencia completa viene incluida en el
% fichero fdl-1.3.tex

% Copyright (C) 2009 Pablo Recio Quijano 
Es usual en un PFC referenciar que software has usado para la
realización del mismo. Aprovecharé este apartado para que conozcas
alguna herramienta que puede serte de ayuda para realizar tus
documentos en \LaTeX{}

\section*{Emacs + Auc\TeX}

Emacs es uno de los programas de edición más usados por
desarrolladores de software, ya que es bastante versatil admitiendo
gran cantidad de ``plugins'' o extensiones que permiten ampliar aun
más sus funcionalidades.\\

Uno de estos plugins es Auc\TeX \cite{pdf:auctex}, el cual incluye
rutas para ciertos comandos, resaltado de sintaxis, previsualización
del documento, menú matemático en el cual podemos acceder e insertar
la gran mayoria de los símbolos matemáticos, para no tener que
memorizarlos. Podemos ver un ejemplo de Emacs + Auc\TeX en la figura
\ref{auctex}

\figura{Auctex.png}{scale=0.6}{Emacs + Auc\TeX}{auctex}{h}

Por ejemplo, para cerrar un entorno $\backslash$\texttt{begin()}, con su
respectivo $\backslash$\texttt{end()}, utilizaremos el atajo
\comando{C-c M-]}, para añadir un $\backslash$\texttt{item}, tenemos
el atajo \comando{C-c C-j}, y así unos cuantos, que una vez que nos
habituamos a ellos, son bastante cómodos.\\

Además, es bastante configurable, con indentado automático, corrector
ortográfico y demás. El fichero adjunto a este documento,
\emph{conf\_emacs} incluye una configuración con varias de estas
opciones.

\section*{Doxygen}

Realmente, \programa{Doxygen} \cite{website:doxygen} no es una herramienta
que vayamos a utilizar para realizar documentos \LaTeX{}
directmaente. Sin embargo, para la documentación de código si es
bastante util.\\

Esta herramienta realiza una documentación automática de código
fuente. Es decir, para nuestro PFC, podemos utilizar para generar la
documentación de las APIs de nuestras librerias y demás. Puede generar
esta documentación en varios formatos, y entre ellos, \LaTeX, de forma
que podemos utilizar ese código generado en nuestra memoria de forma
automática.

\section*{GNU Make}

\programa{GNU Make} es el programa de recompilación y de control de
dependencias por excelencia. Se puede utilizar para compilar proyectos
software en diversos códigos, o como en el caso de este documento,
para compilar documentos \LaTeX{} con diversas opciones.\\

Para más información \cite{pdf:make}

\section*{Dia}

\programa{Dia} es un editor de gráficos vectoriales el cual incluye
distintas plantillas para distintos tipos de gráficos, como pueden ser
UML, ERe, diagramas de flujo, esquemas Cisco de red y un larguísimo
etcétera. Podemos ver el interfaz en la figura \ref{dia}

\figura{dia.png}{scale=0.6}{Interfaz de Dia}{dia}{h}

Estos diagramas podemos exportarlos a diversos formatos de imagen
(\texttt{.png}, \texttt{.eps}, ...) o a formato \texttt{.tex}, como
vimos anteriormente.

\addcontentsline{toc}{chapter}{Instalación de \LaTeX}
\chapter*{Instalación de \LaTeX}
% -*-instalacion.tex-*-
% Este fichero es parte de la plantilla LaTeX para
% la realización de Proyectos Final de Carrera, protejido
% bajo los términos de la licencia GFDL.
% Para más información, la licencia completa viene incluida en el
% fichero fdl-1.3.tex

% Copyright (C) 2009 Pablo Recio Quijano 

Veamos que tenemos que hacer para instalar \LaTeX{} con todas sus
capacidades en un sistema basado en Debian, como Ubuntu.
Primero hay que tener en cuenta que \LaTeX{} es relativamente pesado
con respecto a otros compiladores. \\

Nosotros vamos a utilizar la distribución de \LaTeX{} incluida en los
repositorios de Ubuntu llamada \programa{texlive}. Si la buscas en
tu gestor de paquetes, encontrarás infinidad de paquetes aparte
del principal. Existen otras distribuciones como Te\TeX\\

Si instalas solo los básicos, es decir instalas \programa{texlive} y
los programas necesarios para él, no podrás compilar este documento,
ya que faltarian paquetes tales como \programa{supertabular} y
varios. Por eso, si no tienes problema de espacio en el disco duro te
recomiendo que instales el paquete \programa{texlive-full}, que
instala \negrita{todos} los paquetes de \programa{texlive}, incluyendo
documentación en todos los idiomas disponibles. Si buscas no tener
problemas de dependencias, este es tu método.\\

\begin{lstlisting}[style=consola]
  sudo apt-get install texlive-full
\end{lstlisting}

En caso de querer ser un poco más concreto, en principio puedes
trabajar con la más básica (\programa{texlive} y sus dependencias) y
en función de los paquetes que te vayan faltando, los instalas.

% TO-DO: paquetes concretos e instalación en otros sistemas

% This is set up to run with pdflatex.
%---------The file header---------------------------------------------
%---------------------------------------------------------------------
\chapter*{\rlap{GNU Free Documentation License}}
\phantomsection  % so hyperref creates bookmarks
\addcontentsline{toc}{chapter}{GNU Free Documentation License}
%\label{label_fdl}

 \begin{center}

       Version 1.3, 3 November 2008


 Copyright \copyright{} 2000, 2001, 2002, 2007, 2008  Free Software Foundation, Inc.
 
 \bigskip
 
     <http://fsf.org/>
  
 \bigskip
 
 Everyone is permitted to copy and distribute verbatim copies
 of this license document, but changing it is not allowed.
\end{center}


\begin{center}
{\bf\large Preamble}
\end{center}

The purpose of this License is to make a manual, textbook, or other
functional and useful document ``free'' in the sense of freedom: to
assure everyone the effective freedom to copy and redistribute it,
with or without modifying it, either commercially or noncommercially.
Secondarily, this License preserves for the author and publisher a way
to get credit for their work, while not being considered responsible
for modifications made by others.

This License is a kind of ``copyleft'', which means that derivative
works of the document must themselves be free in the same sense.  It
complements the GNU General Public License, which is a copyleft
license designed for free software.

We have designed this License in order to use it for manuals for free
software, because free software needs free documentation: a free
program should come with manuals providing the same freedoms that the
software does.  But this License is not limited to software manuals;
it can be used for any textual work, regardless of subject matter or
whether it is published as a printed book.  We recommend this License
principally for works whose purpose is instruction or reference.


\begin{center}
{\Large\bf 1. APPLICABILITY AND DEFINITIONS\par}
\phantomsection
\addcontentsline{toc}{section}{1. APPLICABILITY AND DEFINITIONS}
\end{center}

This License applies to any manual or other work, in any medium, that
contains a notice placed by the copyright holder saying it can be
distributed under the terms of this License.  Such a notice grants a
world-wide, royalty-free license, unlimited in duration, to use that
work under the conditions stated herein.  The ``\textbf{Document}'', below,
refers to any such manual or work.  Any member of the public is a
licensee, and is addressed as ``\textbf{you}''.  You accept the license if you
copy, modify or distribute the work in a way requiring permission
under copyright law.

A ``\textbf{Modified Version}'' of the Document means any work containing the
Document or a portion of it, either copied verbatim, or with
modifications and/or translated into another language.

A ``\textbf{Secondary Section}'' is a named appendix or a front-matter section of
the Document that deals exclusively with the relationship of the
publishers or authors of the Document to the Document's overall subject
(or to related matters) and contains nothing that could fall directly
within that overall subject.  (Thus, if the Document is in part a
textbook of mathematics, a Secondary Section may not explain any
mathematics.)  The relationship could be a matter of historical
connection with the subject or with related matters, or of legal,
commercial, philosophical, ethical or political position regarding
them.

The ``\textbf{Invariant Sections}'' are certain Secondary Sections whose titles
are designated, as being those of Invariant Sections, in the notice
that says that the Document is released under this License.  If a
section does not fit the above definition of Secondary then it is not
allowed to be designated as Invariant.  The Document may contain zero
Invariant Sections.  If the Document does not identify any Invariant
Sections then there are none.

The ``\textbf{Cover Texts}'' are certain short passages of text that are listed,
as Front-Cover Texts or Back-Cover Texts, in the notice that says that
the Document is released under this License.  A Front-Cover Text may
be at most 5 words, and a Back-Cover Text may be at most 25 words.

A ``\textbf{Transparent}'' copy of the Document means a machine-readable copy,
represented in a format whose specification is available to the
general public, that is suitable for revising the document
straightforwardly with generic text editors or (for images composed of
pixels) generic paint programs or (for drawings) some widely available
drawing editor, and that is suitable for input to text formatters or
for automatic translation to a variety of formats suitable for input
to text formatters.  A copy made in an otherwise Transparent file
format whose markup, or absence of markup, has been arranged to thwart
or discourage subsequent modification by readers is not Transparent.
An image format is not Transparent if used for any substantial amount
of text.  A copy that is not ``Transparent'' is called ``\textbf{Opaque}''.

Examples of suitable formats for Transparent copies include plain
ASCII without markup, Texinfo input format, LaTeX input format, SGML
or XML using a publicly available DTD, and standard-conforming simple
HTML, PostScript or PDF designed for human modification.  Examples of
transparent image formats include PNG, XCF and JPG.  Opaque formats
include proprietary formats that can be read and edited only by
proprietary word processors, SGML or XML for which the DTD and/or
processing tools are not generally available, and the
machine-generated HTML, PostScript or PDF produced by some word
processors for output purposes only.

The ``\textbf{Title Page}'' means, for a printed book, the title page itself,
plus such following pages as are needed to hold, legibly, the material
this License requires to appear in the title page.  For works in
formats which do not have any title page as such, ``Title Page'' means
the text near the most prominent appearance of the work's title,
preceding the beginning of the body of the text.

The ``\textbf{publisher}'' means any person or entity that distributes
copies of the Document to the public.

A section ``\textbf{Entitled XYZ}'' means a named subunit of the Document whose
title either is precisely XYZ or contains XYZ in parentheses following
text that translates XYZ in another language.  (Here XYZ stands for a
specific section name mentioned below, such as ``\textbf{Acknowledgements}'',
``\textbf{Dedications}'', ``\textbf{Endorsements}'', or ``\textbf{History}''.)  
To ``\textbf{Preserve the Title}''
of such a section when you modify the Document means that it remains a
section ``Entitled XYZ'' according to this definition.

The Document may include Warranty Disclaimers next to the notice which
states that this License applies to the Document.  These Warranty
Disclaimers are considered to be included by reference in this
License, but only as regards disclaiming warranties: any other
implication that these Warranty Disclaimers may have is void and has
no effect on the meaning of this License.


\begin{center}
{\Large\bf 2. VERBATIM COPYING\par}
\phantomsection
\addcontentsline{toc}{section}{2. VERBATIM COPYING}
\end{center}

You may copy and distribute the Document in any medium, either
commercially or noncommercially, provided that this License, the
copyright notices, and the license notice saying this License applies
to the Document are reproduced in all copies, and that you add no other
conditions whatsoever to those of this License.  You may not use
technical measures to obstruct or control the reading or further
copying of the copies you make or distribute.  However, you may accept
compensation in exchange for copies.  If you distribute a large enough
number of copies you must also follow the conditions in section~3.

You may also lend copies, under the same conditions stated above, and
you may publicly display copies.


\begin{center}
{\Large\bf 3. COPYING IN QUANTITY\par}
\phantomsection
\addcontentsline{toc}{section}{3. COPYING IN QUANTITY}
\end{center}


If you publish printed copies (or copies in media that commonly have
printed covers) of the Document, numbering more than 100, and the
Document's license notice requires Cover Texts, you must enclose the
copies in covers that carry, clearly and legibly, all these Cover
Texts: Front-Cover Texts on the front cover, and Back-Cover Texts on
the back cover.  Both covers must also clearly and legibly identify
you as the publisher of these copies.  The front cover must present
the full title with all words of the title equally prominent and
visible.  You may add other material on the covers in addition.
Copying with changes limited to the covers, as long as they preserve
the title of the Document and satisfy these conditions, can be treated
as verbatim copying in other respects.

If the required texts for either cover are too voluminous to fit
legibly, you should put the first ones listed (as many as fit
reasonably) on the actual cover, and continue the rest onto adjacent
pages.

If you publish or distribute Opaque copies of the Document numbering
more than 100, you must either include a machine-readable Transparent
copy along with each Opaque copy, or state in or with each Opaque copy
a computer-network location from which the general network-using
public has access to download using public-standard network protocols
a complete Transparent copy of the Document, free of added material.
If you use the latter option, you must take reasonably prudent steps,
when you begin distribution of Opaque copies in quantity, to ensure
that this Transparent copy will remain thus accessible at the stated
location until at least one year after the last time you distribute an
Opaque copy (directly or through your agents or retailers) of that
edition to the public.

It is requested, but not required, that you contact the authors of the
Document well before redistributing any large number of copies, to give
them a chance to provide you with an updated version of the Document.


\begin{center}
{\Large\bf 4. MODIFICATIONS\par}
\phantomsection
\addcontentsline{toc}{section}{4. MODIFICATIONS}
\end{center}

You may copy and distribute a Modified Version of the Document under
the conditions of sections 2 and 3 above, provided that you release
the Modified Version under precisely this License, with the Modified
Version filling the role of the Document, thus licensing distribution
and modification of the Modified Version to whoever possesses a copy
of it.  In addition, you must do these things in the Modified Version:

\begin{itemize}
\item[A.] 
   Use in the Title Page (and on the covers, if any) a title distinct
   from that of the Document, and from those of previous versions
   (which should, if there were any, be listed in the History section
   of the Document).  You may use the same title as a previous version
   if the original publisher of that version gives permission.
   
\item[B.]
   List on the Title Page, as authors, one or more persons or entities
   responsible for authorship of the modifications in the Modified
   Version, together with at least five of the principal authors of the
   Document (all of its principal authors, if it has fewer than five),
   unless they release you from this requirement.
   
\item[C.]
   State on the Title page the name of the publisher of the
   Modified Version, as the publisher.
   
\item[D.]
   Preserve all the copyright notices of the Document.
   
\item[E.]
   Add an appropriate copyright notice for your modifications
   adjacent to the other copyright notices.
   
\item[F.]
   Include, immediately after the copyright notices, a license notice
   giving the public permission to use the Modified Version under the
   terms of this License, in the form shown in the Addendum below.
   
\item[G.]
   Preserve in that license notice the full lists of Invariant Sections
   and required Cover Texts given in the Document's license notice.
   
\item[H.]
   Include an unaltered copy of this License.
   
\item[I.]
   Preserve the section Entitled ``History'', Preserve its Title, and add
   to it an item stating at least the title, year, new authors, and
   publisher of the Modified Version as given on the Title Page.  If
   there is no section Entitled ``History'' in the Document, create one
   stating the title, year, authors, and publisher of the Document as
   given on its Title Page, then add an item describing the Modified
   Version as stated in the previous sentence.
   
\item[J.]
   Preserve the network location, if any, given in the Document for
   public access to a Transparent copy of the Document, and likewise
   the network locations given in the Document for previous versions
   it was based on.  These may be placed in the ``History'' section.
   You may omit a network location for a work that was published at
   least four years before the Document itself, or if the original
   publisher of the version it refers to gives permission.
   
\item[K.]
   For any section Entitled ``Acknowledgements'' or ``Dedications'',
   Preserve the Title of the section, and preserve in the section all
   the substance and tone of each of the contributor acknowledgements
   and/or dedications given therein.
   
\item[L.]
   Preserve all the Invariant Sections of the Document,
   unaltered in their text and in their titles.  Section numbers
   or the equivalent are not considered part of the section titles.
   
\item[M.]
   Delete any section Entitled ``Endorsements''.  Such a section
   may not be included in the Modified Version.
   
\item[N.]
   Do not retitle any existing section to be Entitled ``Endorsements''
   or to conflict in title with any Invariant Section.
   
\item[O.]
   Preserve any Warranty Disclaimers.
\end{itemize}

If the Modified Version includes new front-matter sections or
appendices that qualify as Secondary Sections and contain no material
copied from the Document, you may at your option designate some or all
of these sections as invariant.  To do this, add their titles to the
list of Invariant Sections in the Modified Version's license notice.
These titles must be distinct from any other section titles.

You may add a section Entitled ``Endorsements'', provided it contains
nothing but endorsements of your Modified Version by various
parties---for example, statements of peer review or that the text has
been approved by an organization as the authoritative definition of a
standard.

You may add a passage of up to five words as a Front-Cover Text, and a
passage of up to 25 words as a Back-Cover Text, to the end of the list
of Cover Texts in the Modified Version.  Only one passage of
Front-Cover Text and one of Back-Cover Text may be added by (or
through arrangements made by) any one entity.  If the Document already
includes a cover text for the same cover, previously added by you or
by arrangement made by the same entity you are acting on behalf of,
you may not add another; but you may replace the old one, on explicit
permission from the previous publisher that added the old one.

The author(s) and publisher(s) of the Document do not by this License
give permission to use their names for publicity for or to assert or
imply endorsement of any Modified Version.


\begin{center}
{\Large\bf 5. COMBINING DOCUMENTS\par}
\phantomsection
\addcontentsline{toc}{section}{5. COMBINING DOCUMENTS}
\end{center}


You may combine the Document with other documents released under this
License, under the terms defined in section~4 above for modified
versions, provided that you include in the combination all of the
Invariant Sections of all of the original documents, unmodified, and
list them all as Invariant Sections of your combined work in its
license notice, and that you preserve all their Warranty Disclaimers.

The combined work need only contain one copy of this License, and
multiple identical Invariant Sections may be replaced with a single
copy.  If there are multiple Invariant Sections with the same name but
different contents, make the title of each such section unique by
adding at the end of it, in parentheses, the name of the original
author or publisher of that section if known, or else a unique number.
Make the same adjustment to the section titles in the list of
Invariant Sections in the license notice of the combined work.

In the combination, you must combine any sections Entitled ``History''
in the various original documents, forming one section Entitled
``History''; likewise combine any sections Entitled ``Acknowledgements'',
and any sections Entitled ``Dedications''.  You must delete all sections
Entitled ``Endorsements''.

\begin{center}
{\Large\bf 6. COLLECTIONS OF DOCUMENTS\par}
\phantomsection
\addcontentsline{toc}{section}{6. COLLECTIONS OF DOCUMENTS}
\end{center}

You may make a collection consisting of the Document and other documents
released under this License, and replace the individual copies of this
License in the various documents with a single copy that is included in
the collection, provided that you follow the rules of this License for
verbatim copying of each of the documents in all other respects.

You may extract a single document from such a collection, and distribute
it individually under this License, provided you insert a copy of this
License into the extracted document, and follow this License in all
other respects regarding verbatim copying of that document.


\begin{center}
{\Large\bf 7. AGGREGATION WITH INDEPENDENT WORKS\par}
\phantomsection
\addcontentsline{toc}{section}{7. AGGREGATION WITH INDEPENDENT WORKS}
\end{center}


A compilation of the Document or its derivatives with other separate
and independent documents or works, in or on a volume of a storage or
distribution medium, is called an ``aggregate'' if the copyright
resulting from the compilation is not used to limit the legal rights
of the compilation's users beyond what the individual works permit.
When the Document is included in an aggregate, this License does not
apply to the other works in the aggregate which are not themselves
derivative works of the Document.

If the Cover Text requirement of section~3 is applicable to these
copies of the Document, then if the Document is less than one half of
the entire aggregate, the Document's Cover Texts may be placed on
covers that bracket the Document within the aggregate, or the
electronic equivalent of covers if the Document is in electronic form.
Otherwise they must appear on printed covers that bracket the whole
aggregate.


\begin{center}
{\Large\bf 8. TRANSLATION\par}
\phantomsection
\addcontentsline{toc}{section}{8. TRANSLATION}
\end{center}


Translation is considered a kind of modification, so you may
distribute translations of the Document under the terms of section~4.
Replacing Invariant Sections with translations requires special
permission from their copyright holders, but you may include
translations of some or all Invariant Sections in addition to the
original versions of these Invariant Sections.  You may include a
translation of this License, and all the license notices in the
Document, and any Warranty Disclaimers, provided that you also include
the original English version of this License and the original versions
of those notices and disclaimers.  In case of a disagreement between
the translation and the original version of this License or a notice
or disclaimer, the original version will prevail.

If a section in the Document is Entitled ``Acknowledgements'',
``Dedications'', or ``History'', the requirement (section~4) to Preserve
its Title (section~1) will typically require changing the actual
title.


\begin{center}
{\Large\bf 9. TERMINATION\par}
\phantomsection
\addcontentsline{toc}{section}{9. TERMINATION}
\end{center}


You may not copy, modify, sublicense, or distribute the Document
except as expressly provided under this License.  Any attempt
otherwise to copy, modify, sublicense, or distribute it is void, and
will automatically terminate your rights under this License.

However, if you cease all violation of this License, then your license
from a particular copyright holder is reinstated (a) provisionally,
unless and until the copyright holder explicitly and finally
terminates your license, and (b) permanently, if the copyright holder
fails to notify you of the violation by some reasonable means prior to
60 days after the cessation.

Moreover, your license from a particular copyright holder is
reinstated permanently if the copyright holder notifies you of the
violation by some reasonable means, this is the first time you have
received notice of violation of this License (for any work) from that
copyright holder, and you cure the violation prior to 30 days after
your receipt of the notice.

Termination of your rights under this section does not terminate the
licenses of parties who have received copies or rights from you under
this License.  If your rights have been terminated and not permanently
reinstated, receipt of a copy of some or all of the same material does
not give you any rights to use it.


\begin{center}
{\Large\bf 10. FUTURE REVISIONS OF THIS LICENSE\par}
\phantomsection
\addcontentsline{toc}{section}{10. FUTURE REVISIONS OF THIS LICENSE}
\end{center}


The Free Software Foundation may publish new, revised versions
of the GNU Free Documentation License from time to time.  Such new
versions will be similar in spirit to the present version, but may
differ in detail to address new problems or concerns.  See
http://www.gnu.org/copyleft/.

Each version of the License is given a distinguishing version number.
If the Document specifies that a particular numbered version of this
License ``or any later version'' applies to it, you have the option of
following the terms and conditions either of that specified version or
of any later version that has been published (not as a draft) by the
Free Software Foundation.  If the Document does not specify a version
number of this License, you may choose any version ever published (not
as a draft) by the Free Software Foundation.  If the Document
specifies that a proxy can decide which future versions of this
License can be used, that proxy's public statement of acceptance of a
version permanently authorizes you to choose that version for the
Document.


\begin{center}
{\Large\bf 11. RELICENSING\par}
\phantomsection
\addcontentsline{toc}{section}{11. RELICENSING}
\end{center}


``Massive Multiauthor Collaboration Site'' (or ``MMC Site'') means any
World Wide Web server that publishes copyrightable works and also
provides prominent facilities for anybody to edit those works.  A
public wiki that anybody can edit is an example of such a server.  A
``Massive Multiauthor Collaboration'' (or ``MMC'') contained in the
site means any set of copyrightable works thus published on the MMC
site.

``CC-BY-SA'' means the Creative Commons Attribution-Share Alike 3.0
license published by Creative Commons Corporation, a not-for-profit
corporation with a principal place of business in San Francisco,
California, as well as future copyleft versions of that license
published by that same organization.

``Incorporate'' means to publish or republish a Document, in whole or
in part, as part of another Document.

An MMC is ``eligible for relicensing'' if it is licensed under this
License, and if all works that were first published under this License
somewhere other than this MMC, and subsequently incorporated in whole
or in part into the MMC, (1) had no cover texts or invariant sections,
and (2) were thus incorporated prior to November 1, 2008.

The operator of an MMC Site may republish an MMC contained in the site
under CC-BY-SA on the same site at any time before August 1, 2009,
provided the MMC is eligible for relicensing.


\begin{center}
{\Large\bf ADDENDUM: How to use this License for your documents\par}
\phantomsection
\addcontentsline{toc}{section}{ADDENDUM: How to use this License for your documents}
\end{center}

To use this License in a document you have written, include a copy of
the License in the document and put the following copyright and
license notices just after the title page:

\bigskip
\begin{quote}
    Copyright \copyright{}  YEAR  YOUR NAME.
    Permission is granted to copy, distribute and/or modify this document
    under the terms of the GNU Free Documentation License, Version 1.3
    or any later version published by the Free Software Foundation;
    with no Invariant Sections, no Front-Cover Texts, and no Back-Cover Texts.
    A copy of the license is included in the section entitled ``GNU
    Free Documentation License''.
\end{quote}
\bigskip
    
If you have Invariant Sections, Front-Cover Texts and Back-Cover Texts,
replace the ``with \dots\ Texts.'' line with this:

\bigskip
\begin{quote}
    with the Invariant Sections being LIST THEIR TITLES, with the
    Front-Cover Texts being LIST, and with the Back-Cover Texts being LIST.
\end{quote}
\bigskip
    
If you have Invariant Sections without Cover Texts, or some other
combination of the three, merge those two alternatives to suit the
situation.

If your document contains nontrivial examples of program code, we
recommend releasing these examples in parallel under your choice of
free software license, such as the GNU General Public License,
to permit their use in free software.

%---------------------------------------------------------------------


\end{document}
